\documentclass[10pt,a4paper]{article}

\usepackage[margin=16mm,headheight=15pt]{geometry}
\usepackage{fontspec}
\usepackage{kotex}
\usepackage{xcolor}
\usepackage{array}
\usepackage{booktabs}
\usepackage{longtable}
\usepackage{tabularx}
\usepackage{enumitem}
\usepackage{amsmath,amssymb,bm}
\usepackage[most]{tcolorbox}
\usepackage{hyperref}
\usepackage{fancyhdr}
\usepackage{listings}

\setmainfont{Noto Serif CJK KR}
\setsansfont{Noto Sans CJK KR}
\setmonofont{DejaVu Sans Mono}

\hypersetup{
  colorlinks=true,
  linkcolor=blue!50!black,
  urlcolor=blue!50!black,
  pdftitle={Topology-Distilled, Error-Triggered Global--Local Wind Dynamics Implementation Checklist},
  pdfauthor={Wind3DGS}
}

\definecolor{TodoGray}{RGB}{105,105,105}
\definecolor{ProgressBlue}{RGB}{35,90,170}
\definecolor{DoneGreen}{RGB}{35,130,80}
\definecolor{BlockRed}{RGB}{170,55,55}
\definecolor{KillRed}{RGB}{125,20,25}
\definecolor{LightBlue}{RGB}{235,243,255}
\definecolor{LightGray}{RGB}{248,248,248}
\definecolor{LightRed}{RGB}{255,239,239}
\definecolor{LightGreen}{RGB}{238,250,242}
\definecolor{LightOrange}{RGB}{255,247,233}

\newcommand{\todo}{\textcolor{TodoGray}{\(\square\)}}
\newcommand{\progress}{\textcolor{ProgressBlue}{[{\raisebox{0.1ex}{\(\sim\)}}]}}
\newcommand{\done}{\textcolor{DoneGreen}{\(\boxtimes\)}}
\newcommand{\blocked}{\textcolor{BlockRed}{[!]}}
\newcommand{\killmark}{\textcolor{KillRed}{[K]}}
\newcommand{\code}[1]{\texttt{#1}}
\newcommand{\artifact}[1]{\textcolor{ProgressBlue}{\nolinkurl{#1}}}
\newcolumntype{L}[1]{>{\raggedright\arraybackslash}p{#1}}
\newcolumntype{Y}{>{\raggedright\arraybackslash}X}

\setlist[itemize]{leftmargin=1.45em,itemsep=0.16em,topsep=0.28em}
\setlist[itemize,2]{label=--,leftmargin=1.8em}
\setlist[enumerate]{leftmargin=1.65em,itemsep=0.2em,topsep=0.3em}
\setlength{\emergencystretch}{2.5em}
\renewcommand{\arraystretch}{1.2}
\sloppy

\lstset{
  basicstyle=\ttfamily\small,
  frame=single,
  breaklines=true,
  columns=fullflexible,
  backgroundcolor=\color{LightGray},
  rulecolor=\color{black!20},
  xleftmargin=0.4em,
  xrightmargin=0.4em,
  aboveskip=0.5em,
  belowskip=0.5em
}

\pagestyle{fancy}
\fancyhf{}
\lhead{Topology-Distilled Error-Triggered Wind Dynamics}
\rhead{Implementation Checklist}
\cfoot{\thepage}

\title{\vspace{-1.7em}\textbf{Topology-Distilled, Error-Triggered\\Global--Local Wind Dynamics}\\[0.4em]\Large 구현 체크리스트}
\author{Wind3DGS}
\date{Created: 2026-08-13 \quad Status baseline: new method reset}

\begin{document}
\maketitle

\begin{tcolorbox}[
  colback=LightBlue,
  colframe=blue!45!black,
  title=고정 Runtime Pipeline,
  fonttitle=\bfseries
]
\begin{center}
\code{static 3DGS + setup metadata $\rightarrow$ versioned object package}
\(\rightarrow\)
\code{current surface + full-anchor analytic aero}
\(\rightarrow\)
\code{Global predictor}
\(\rightarrow\)
\code{fixed-$K$ cached-compliance selector (Oracle upper bound)}
\(\rightarrow\)
\code{active Local physical base}
\\[0.25em]
\(\rightarrow\)
\code{joint overlap/conservation/complement assembly}
\(\rightarrow\)
\code{single Local physical solve}
\(\rightarrow\)
\code{Global corrector}
\(\rightarrow\)
\code{Gaussian transport/render}
\end{center}

\textbf{V0 core}는 full-anchor aero, patch-local fixed-superset basis,
fixed-cardinality selector, physics-only Local, 보존 조립/complement/corrector와 affine transport다.
다음 세 모듈은 core dependency가 아니라 독립적인 go/no-go 후보다.
\textbf{H}: aerodynamic hyper-reduction,
\textbf{F}: learned missing-force expert,
\textbf{G}: learned future-benefit gate이며, hardware cost-aware selection은 profiler가 필요성을 보일 때의 후속 확장이다.
각 후보는 대응 profiling 또는 Oracle upper bound를 통과한 경우에만 활성화한다.

``Error-triggered''는 certified posterior estimator를 뜻하지 않는다.
Learned G의 존재가 아니라 cached-compliance가 held-out object/gust에서 Oracle benefit과의
rank correlation, regret, activation delay/high-benefit miss와 actual-skip quality--latency를 통과하는지가
최종 제목의 evidence gate다. 통과하지 못하면 결과에 따라
\emph{Compliance-Guided} 또는 \emph{Selective Complementary}로 낮추며,
every-frame score가 residual이 아니므로 \emph{Residual-Guided}라고 부르지 않는다.

이 문서는 새 연구 방향의 구현 상태를 관리하는 source of truth다.
방법 정의와 수식의 source of truth는
\path{ideas/3dgs_topology_distilled_error_triggered_wind_dynamics_2026-08-13.tex}다.
이전 proxy/SH-residual 체크리스트의 완료 상태는 승계하지 않는다.
\end{tcolorbox}

\begin{tcolorbox}[colback=LightRed,colframe=KillRed,title=가장 중요한 실행 순서 계약,fonttitle=\bfseries]
논문 설명의 Module 번호는 이해를 위한 논리 번호다. 실제 구현은 다음 dependency를 따른다.
\begin{center}
\code{1 $\rightarrow$ 2 $\rightarrow$ 3 $\rightarrow$ 4 $\rightarrow$ 5 $\rightarrow$ 6 $\rightarrow$ 7
$\rightarrow$ 10 $\rightarrow$ 9 $\rightarrow$ 8 $\rightarrow$ 11 $\rightarrow$ 12 $\rightarrow$ 13 $\rightarrow$ 14}
\end{center}
Module 7의 analytic Local base aero와, 조건부 F를 사용할 때의 missing aero/missing structural을 channel별로 유지한다.
존재하는 channel을 Module 10에서 먼저 조립하고 conservation과
\(\Phi^\top f=0\)를 하나의 joint KKT로 만족시킨다.
Module 9는 contract/generalized projection만 검사하며 force를 다시 순차 projection하지 않는다.
Module 8에서만 channel을 합쳐 단일 Local system을 한 번 적분한다.
Patch별로 적분한 뒤 위치를 평균하는 구현은 금지한다.
\end{tcolorbox}

\tableofcontents
\newpage

\section{문서 사용 규칙}

\subsection{상태 표기}

\begin{center}
\begin{tabular}{@{}ll@{}}
\toprule
상태 & 의미 \\
\midrule
\todo & 아직 시작하지 않음 \\
\progress & 구현 또는 검증 진행 중 \\
\done & 자동 테스트와 완료 기준을 모두 통과함 \\
\blocked & 설계 결정, 외부 자산 또는 환경 문제로 막힘 \\
\killmark & kill gate 발생; 원인 해결 전 후속 milestone 중단 \\
\bottomrule
\end{tabular}
\end{center}

\begin{itemize}
  \item 구현 코드만 존재하는 상태는 \done이 아니다. 재현 명령, 자동 검증, 산출물까지 확인되어야 한다.
  \item Manual inspection이 남아 있으면 최대 \progress로 둔다.
  \item \blocked 항목에는 반드시 \textbf{현재 상태, 원인, 다음 결정}을 기록한다.
  \item \killmark가 발생하면 성능 수치를 완화하여 통과시키지 않고 upstream 설계를 고친다.
\end{itemize}

\subsection{Canonical 수식과 section 참조 규칙}

수식의 정의는 canonical idea sketch에만 둔다. 이 체크리스트와 code artifact는
PDF에서 재배치될 수 있는 식 번호나 ``식 17--63'' 같은 숫자 범위를 복사하지 않고,
sketch의 안정적인 LaTeX label 또는 section title을 기록한다.
예를 들어 relation confidence는 \code{eq:relation-validity-confidence}와
\code{eq:relation-absolute-confidence}, structural freeze는
\code{eq:canonical-structural-force-ledger},
\code{eq:klin-projective-hard-gate},
\code{eq:v0-structural-freeze-signature}, support/reaction은
\code{eq:support-restricted-global-routing},
\code{eq:hard-total-reaction-residual},
\code{eq:hard-reaction-rom-residual-decomposition}을 source로 삼는다.
외부 문서의 label이라 이 PDF에서 \code{\textbackslash eqref}로 재해석하지 않고
manifest의 \code{canonical\_source\_path}, \code{canonical\_source\_commit},
\code{canonical\_labels[]}로 추적한다. Label이 바뀌면 checklist와 schema reference test가 함께 실패해야 한다.

\subsection{자동 목표 실행 프로토콜}

권장 요청 형식은 다음과 같다.

\begin{lstlisting}
[Wind3DGS | code | TD05] 목표: TD05 완료까지 구현-검증 루프 실행.
중지 조건: 동일 error signature 3회, kill gate, 사용자 판단 필요,
          검증 환경/데이터 부재, 승인 필요한 외부 변경.
\end{lstlisting}

\begin{enumerate}
  \item 각 루프는 \code{구현 -> 검증 -> 수정 -> 재검증} 순서로 진행한다.
  \item 오류는 실행 명령, 실패 check, 핵심 traceback/message를 묶은 error signature로 추적한다.
  \item 같은 signature를 고치는 시도는 최대 3회다. 이후에는 실패 근거와 선택지를 보고한다.
  \item Dataset 선택, 라이선스, 장시간 GPU 학습, system package 설치, 완료 기준 변경,
  외부 서비스, destructive operation은 사용자 판단을 받는다.
  \item 의미 있는 루프가 끝나면 해당 work folder의 \path{sessions/}에
  목표, 변경 파일, 검증 명령, 결과, blocker, 다음 액션을 기록한다.
\end{enumerate}

\subsection{Experiment와 코드 naming convention}

기존 \code{M\#\#} 실험과 충돌하지 않도록 새 방법은 \code{TD\#\#} tag를 사용한다.

\begin{tcolorbox}[colback=LightGray,colframe=black!20,title=권장 경로,fonttitle=\bfseries]
\begin{itemize}
  \item Experiment: \path{experiments/TD##_short_name/}
  \item Reusable code: \path{code/wind3dgs/}
  \item Idea/checklist: \path{ideas/}
  \item Training output/checkpoint: experiment output 또는 명시된 artifact store
\end{itemize}
\end{tcolorbox}

Reusable code는 semantic package로 나눈다.

\begin{lstlisting}
code/wind3dgs/
  contracts/      # schemas, frames, units, force ledger
  teacher/        # teacher conversion and label generation
  topology/       # anchor graph and operator distillation
  aero/           # full-anchor current-surface law; optional H backend
  reduced/        # Global/Local bases and integrators
  local/          # patch proposals, assembly, complement, feedback
  learning/       # topology; optional F missing-force and G gate models
  runtime/        # typed 14-module engine
  transport/      # affine anchor-to-Gaussian transport
  evaluation/     # metrics, baselines, profiling
\end{lstlisting}

\section{Scope reset와 현재 상태}

\subsection{이번 체크리스트의 runtime contract}

\begin{longtable}{L{0.23\linewidth}L{0.18\linewidth}L{0.51\linewidth}}
\toprule
항목 & Runtime 허용 & 계약 \\
\midrule
\endfirsthead
\toprule
항목 & Runtime 허용 & 계약 \\
\midrule
\endhead
Canonical static 3DGS & 허용 & Target의 기본 입력이다. \\
Prescribed spatial wind query & 허용 & 유체 상태를 푸는 것이 아니라 control/input으로 조회한다. \\
Sparse point/edge/local-relation scaffold & 허용 & Local tangent frame과 area weight를 가진 non-watertight graph다. Global face connectivity, manifold 또는 watertight surface를 요구하지 않는다. \\
Versioned reduced object package & 허용 & Static GS, metric scale, user material과 user attachment layout을 받아 setup에서 만들고 frame마다 재구성하지 않는다. \\
Explicit target triangle-mesh reconstruction & Mainline 금지 & 별도 end-to-end baseline/fallback에서만 허용하고 결과를 구분한다. \\
Full FEM/MPM/CFD/LBM state & 금지 & Teacher와 accuracy upper bound에서만 사용한다. \\
Frame-wise Gaussian optimization & 금지 & Final anchor state로 affine transport한 뒤 standard renderer를 사용한다. \\
RSH coefficient cache & Mainline 제외 & TD06/E2의 representation ablation으로만 평가한다. \\
Aerodynamic hyper-reduction (H) & 조건부 & Full-anchor profiler가 병목을 보이고 generalized-load fidelity와 wall-clock가 함께 개선될 때만 사용한다. \\
Missing-force expert (F) & 조건부 & Physics-only Local보다 Oracle/always-on 평가에서 반복 가능한 이득이 있을 때만 사용한다. \\
Learned future-benefit gate (G) & 조건부 & Oracle--cached-compliance gap이 존재하고, 같은 atomic grouping/\(K_A\)에서 cached-compliance보다 좋은 measured-latency Pareto를 만들며 Oracle upper bound에 접근할 때만 사용한다. Cost-aware variant는 profiler 이후다. \\
\bottomrule
\end{longtable}

\begin{tcolorbox}[colback=LightBlue,colframe=blue!45!black,title=Setup과 frame runtime의 경계,fonttitle=\bfseries]
Static appearance에서 metric material이나 attachment layout을 자동 식별한다고 가정하지 않는다.
\code{SetupInput = (G0, metric scale 또는 nondimensional preset, user material metadata, user attachment layout)}이고,
그 결과가 versioned \code{ObjectPackage}다.
Attachment layout 또는 basis를 바꾸는 material 값을 수정하면 package/operator/basis를 다시 만들거나 refactorize한다.
첫 논문에서 frame runtime control은 prescribed wind와, 이미 고정된 support DOF의 handle motion으로 제한한다.
\end{tcolorbox}

\subsection{기존 구현의 재사용 후보}

이전 체크리스트의 \done/\progress 상태를 자동으로 옮기지 않는다.
다음 항목은 새 schema와 E0 contract를 다시 통과하기 전까지 재사용 후보일 뿐이다.

\begin{longtable}{L{0.06\linewidth}L{0.31\linewidth}L{0.55\linewidth}}
\toprule
상태 & 후보 & 새 방법에서 필요한 재검증 \\
\midrule
\endfirsthead
\toprule
상태 & 후보 & 새 방법에서 필요한 재검증 \\
\midrule
\endhead
\progress & Static 3DGS PLY loader/camera loader & 새 \code{CanonicalGaussianAsset} schema, 좌표계, dtype, covariance convention 통과. \\
\progress & Synthetic cloth/leaf fixtures와 debug viewer & Oracle strip/flag fixtures로 재사용하되 물리 정확도 완료로 간주하지 않음. \\
\progress & Position/covariance transport prototype & Final total anchor state에서 단 한 번 수행하는 affine MLS, SPD, normal consistency 재검증. \\
\progress & Procedural wind preview & Wind query fixture와 qualitative smoke test로만 사용; teacher 또는 physical solver가 아님. \\
\blocked & \code{gsplat 1.5.3} on GTX 1080 Ti & Compute Capability 6.1 환경의 기존 JIT incompatibility. 논문 rendering은 compatible backend/GPU가 필요. \\
\bottomrule
\end{longtable}

\subsection{현재 상태 baseline}

\begin{itemize}
  \item[\done] 새 아이디어 스케치, 전용 참고문헌, PDF와 bundle이 존재한다.
  \item[\done] 이전 아이디어 문서와 구현 체크리스트는 backup에 보존했다.
  \item[\done] 새 구현 체크리스트와 milestone dependency를 확정한다.
  \item[\todo] 새 방법의 typed schema와 deterministic contract를 코드로 구현한다.
  \item[\todo] Oracle-topology Global MVP부터 새 completion evidence를 생성한다.
\end{itemize}

\section{고정 입력, 출력, 물리량 owner 계약}

\subsection{Runtime 입력과 출력}

\begin{lstlisting}
SetupInput
  canonical_gaussians: G0
  metric_scale_or_nondimensional_preset
  user_material_metadata
  user_attachment_layout

SetupOutput
  object_package: versioned scaffold/operators/patch_bases/binding
  optional_hyper_reduction_backend

FrameRuntimeInput
  object_package
  wind_query: W(x, t)
  prescribed_handle_motion_on_fixed_support_dofs
  dt

FrameOutput
  deformed_gaussians: Gt
  global_state: q, qdot
  local_state: z, zdot
  patch_state: active_set, decay_set, dormant_set
  diagnostics
  optional_rendered_frame
\end{lstlisting}

\subsection{Force와 state owner ledger}

\begin{longtable}{L{0.22\linewidth}L{0.28\linewidth}L{0.42\linewidth}}
\toprule
물리량 & 유일한 owner & 구현 계약 \\
\midrule
\endfirsthead
\toprule
물리량 & 유일한 owner & 구현 계약 \\
\midrule
\endhead
Mass/inertia & Reduced Global/Local solver & Network가 inertia 또는 acceleration 전체를 다시 예측하지 않는다. \\
Restoring force & \(K_{gg},K_{\ell\ell}\), 실제 nonzero cross block & Selector/G와 무관하게 항상 존재한다. \\
Damping & \(C_{gg},C_{\ell\ell}\), 실제 nonzero cross block & Selector-off에서도 Local energy를 물리적으로 감쇠한다. \\
Gravity & Explicit external force & Teacher residual에서 정확히 한 번 제거한다. \\
Attachment & Constraint 또는 compliant force 중 하나 & 같은 support를 projection과 force로 이중 적용하지 않는다. \\
Base aerodynamic force & Current-surface analytic law & Relative wind, current normal/area를 사용한다. \\
Missing force & 조건부 F residual network & F가 채택된 경우에만 base model이 설명한 힘을 제거한 defect를 예측한다. V0에서는 0이다. \\
Patch overlap & Constrained force assembly & Network나 displacement averaging으로 seam을 숨기지 않는다. \\
Global complement & Complement projector/basis & Local motion/force의 Global double counting만 제거한다. \\
Global feedback & Corrected-minus-predictor delta & Predictor가 이미 사용한 base/cross term을 다시 넣지 않는다. \\
Gaussian deformation & Final total anchor state & Global/Local이 Gaussian을 각각 두 번 transport하지 않는다. \\
\bottomrule
\end{longtable}

\begin{tcolorbox}[colback=LightOrange,colframe=orange!60!black,title=Selector가 끄는 것과 끄지 않는 것,fonttitle=\bfseries]
V0 fixed-\(K\) selector와 조건부 G gate는 새 Local excitation을 제어하고, F를 채택했을 때만 expensive missing-force expert 실행도 제어한다.
Global/Local mass, stiffness, damping, gravity, attachment, base aero는 끄지 않는다.
Active에서 Decay로 이동할 때 \(z,\dot z\)를 freeze/reset하지 않고 physical decay solve를 유지한다.
\end{tcolorbox}

\section{P0 implementation kill gates}

아래 항목은 단순 checklist가 아니라 후속 milestone 진행을 금지하는 gate다.

\begin{longtable}{L{0.06\linewidth}L{0.45\linewidth}L{0.41\linewidth}}
\toprule
ID & 실패 조건 & 필요한 조치 \\
\midrule
\endfirsthead
\toprule
ID & 실패 조건 & 필요한 조치 \\
\midrule
\endhead
K00 & 좌표계, 단위, normal/front--back, index convention이 artifact마다 다르다. & TD01 schema와 converter를 고치기 전 dataset/package 생성 금지. \\
K01 & Patch별 \(\|\Phi^\top M\Psi_r\|\)가 tolerance를 넘거나 전역 QR/SVD가 patch identity를 섞는다. & Patch-local complement/normalization과 fixed coordinate owner를 고치기 전 Local solve 금지. \\
K02 & Reachable active/decay block의 \(M\succ0\), typed \(K_{\mathrm{lin}}^0,C_{\mathrm{lin}}^0,K^0,C^0\succeq0\), rank/conditioning 또는 inter-patch symmetry/reciprocity가 깨진다. & Runtime solve 금지; operator/basis/group construction 수정. \\
K03 & Gravity, attachment, aero, projective target/restoring, residual에 owner가 둘 이상이다. & Force/target ledger 수정 후 label과 rollout 재생성. \\
K04 & Selector-off/zero-wind에서 rest 복귀가 실패하거나 Local energy가 증가한다. & Activation, damping, integrator, decay state 수정. \\
K05 & Patch별 integrate-then-average 코드가 있다. & Force assembly 후 단일 Local solve로 구조 변경. \\
K06 & Assembly 후 Global leakage, overlap seam, support-route/joint-Schur per-RHS feasibility,
adjustment/reaction amplification/generalized work 또는 free/hard/compliant reaction ledger가 tolerance를 넘는다. & Joint solve, support/route 또는 basis/patch를 수정한다. Ridge, row 삭제나 sequential projection으로 통과 처리하지 않는다. \\
K07 & Corrector가 corrected-minus-predictor delta 외 base/cross force를 재사용한다. & Consumed-force ledger와 delta 정의 수정. \\
K08 & Inactive 전환 시 \(z,\dot z\)를 즉시 freeze/reset한다. & Active/Decay/Dormant state와 energy threshold 구현. \\
K09 & 3\(\times\)2 tangent area-normal, 별도 3\(\times\)3 full shell map, rigid/affine reproduction 또는 covariance SPD가 실패한다. & Aero/Gaussian output 금지; surface/transport map 계약 수정. \\
K10 & Oracle Local이 동일 p95 latency의 Global rank sweep보다 velocity/trajectory/spectrum을 개선하지 못한다. & Local branch 제거 또는 Global/Local basis 재설계. Coordinate/rank match는 보조 진단으로만 사용. \\
K11 & 조건부 F always-on missing force가 physics-only Local을 개선하지 못한다. & F를 제거하고 pure complementary physics를 유지한다. \\
K12 & 조건부 G learned gate가 fixed-\(K\) cached-compliance selector보다 quality--latency Pareto를 개선하지 못한다. & G/cost model을 제거하고 cached-compliance selector를 유지한다. \\
K13 & Same-anchor sparse solver 대비 개발용 내부 기준인 dynamics-only 약 2배 이득도,
동일 p95 latency의 velocity/trajectory/spectrum 개선도, multi-object throughput 이득도 없다.
& SYS selective-compute/quality--latency headline을 철회한다. 이 실패만으로 K10 Local necessity나 연구-level Gate A를 기각하지 않는다. \\
K14 & Full-anchor aero가 병목이 아니거나 조건부 H가 generalized-load fidelity와 wall-clock를 함께 개선하지 못한다. & H/cubature를 runtime과 main contribution에서 제거한다. \\
\addlinespace
K15 & 고정 physical budget에서 \(N_G\) 증가와 함께 structural dynamics 비용이 증가한다. & Resolution-decoupling system evidence를 철회하고 binding/runtime breakdown을 재검토한다. \\
\addlinespace
K16 & Selector, joint assembly/Schur, active/decay gather, factorization과 solve 때문에 adaptive p95/p99가 always-on Local보다 나쁘다. & Backend/cache/dispatcher를 재검토하고 실패하면 always-on으로 후퇴하며 actual-skip headline을 철회한다. \\
\addlinespace
K17 & Incremental/PCG/independent 근사 backend가 same-set exact direct 대비 linear residual, state 또는 energy tolerance를 넘는다. & 해당 근사만 제거하고 exact coupled backend를 유지한다. \\
\addlinespace
K18 & Stationary/slow/fast/reversing gust에서 \(\rho_{\mathcal S}\simeq1\)이 지속되거나 decay/cache churn 때문에 latency 이득이 사라진다. & Physical decay를 끄거나 state를 reset하지 않고 dispatcher/decay를 재설계한다. 실패하면 always-on으로 후퇴한다. \\
\addlinespace
K19 & Production trace가 full anchor-primal generic KKT, patch 수에 비례하는 duplicate overlap/velocity row, stale factor reuse, sequential fallback 또는 반복적인 support/halo cap 초과를 사용한다. & Unique-anchor support-restricted routing과 동등한 small-Schur/null-space backend로 고친다. Joint 계약을 유지해도 latency 이득이 없으면 actual-skip/SYS headline을 축소한다. \\
\addlinespace
K20 & \code{eq:canonical-structural-force-ledger}의 pure-rigid/strained-state objectivity,
normal-curvature pollution, relative-rate damping, midpoint phase/PSD 또는 one-shot/two-step/converged splitting tolerance가 실패한다. & Track A0에 진입하지 않고 projective target/operator contract를 재설계한다. Asset별 backend 변경이나 backward Euler의 수치 감쇠로 숨기지 않는다. \\
\addlinespace
K21 & Normalized area share는 존재하지만 \code{eq:relation-absolute-confidence}의
calibrated absolute confidence가 낮거나 candidate-count/density/valence/conditioning별 calibration 없이 learned package가 통과한다. & Step 0에서는 schema, synthetic aggregator와 fallback fixture만 검사한다. 정식 calibration acceptance는 Track B0/E1에서 판정하고, 실패하면 support 확대 후 user hint/fallback/reject한다. Confidence를 stiffness에 곱하지 않는다. \\
\addlinespace
K22 & \code{eq:klin-projective-hard-gate}의 energy, stiffness/damping action,
action-tangent symmetry/PSD, modal angle 또는 common-probe closure가 실패한다. & Canonical package는 실패다. Action-validity 실패는 backend를 재설계하고, rest-linear agreement만 실패한 action-basis는 MC1 claim을 바꾼 별도 versioned pivot으로만 검토한다. Learned residual이나 raw \(K_{\mathrm{proj}}^0\) mode 비교로 우회하지 않는다. \\
\addlinespace
K23 & Canonical V0가 \code{eq:fixed-global-sample-jacobian}의 setup-fixed
\(T_s,\Phi_s\) 대신 state-dependent Jacobian을 암묵적으로 만들거나 update/corrector 비용을 숨긴다. & V0를 fixed lifting으로 고치고 nonlinear sample map은 별도 versioned optional backend/profiler로 분리한다. \\
\addlinespace
K24 & Mandatory selector blind-spot fixture에서 cached compliance가 큰 Oracle-benefit region을 반복해서 놓치고,
hybrid/full-counter/조건부 G 대안도 measured-p95 Pareto를 개선하지 못한다. & Error-triggered/adaptive selector claim을 축소하고 always-on Local, merged group 또는 단순 dispatcher로 후퇴한다. \\
\bottomrule
\end{longtable}

연구-level Gate A는 별도의 5--10개 independent-GS representation preflight로
close-layer false edge, area/mass, 초기 8--16 modes/modal subspace를 판정한다.
Gate B는 K10, Gate C는 K13/K16/K18이다.
이 세 gate는 K00--K09, K17/K19와 K20--K24의 구현 invariant를 대체하지 않는다.
Gate A 실패는 MC1 representation claim만 축소하며 Gate B/C는 Oracle 또는
명시적 fallback scaffold에서 계속 판정할 수 있다.
Gate B 실패는 Local/selector claim을 제거하지만 MC1과 Global physics를 기각하지 않는다.
Gate C 실패는 selective-compute/SYS performance headline만 제거하며
Gate A와 Gate B의 accuracy 판정을 뒤집지 않는다.

\section{Milestone dependency와 전체 현황}

\subsection{Step 0 structural/operator-load preflight}

Track A0 또는 대규모 teacher/package 학습 전에 최소 synthetic/oracle strip fixture에서
다음 순서를 고정한다. 최소 fixture와 independent-GS representation probe는 병렬 생성할 수 있지만,
이 순서를 통과하기 전에 full predicted-topology training이나 H/F/G 구현을 시작하지 않는다.

\begin{enumerate}
  \item \(\mathcal Q_{\mathrm{op}}\), metric MLS, area/mass/material/attachment 조립과
  relation validity--anchor absolute-confidence schema/fallback을 typed contract로 고정한다.
  \item Objective normal-curvature의 planar in-plane bending-pollution,
  normal-curvature, shallow-rest와 pure-rigid/strained-state rotation fixture를 실행한다.
  \item Rest-linear/locally refit energy, held-target stiffness/damping action,
  action-tangent symmetry/PSD, modal subspace와 common-probe closure를 판정한다.
  \item One-shot/two-step/converged target fit에서 torque/phase/target-band PSD를 비교하고,
  canonical objective projective target + implicit midpoint를 동결한다.
  Backward Euler와 rest-linear runtime은 ablation으로만 남긴다.
  \item Setup-fixed \(T_s,\Phi_s\) Global lifting과 corrector가 같은 계약을 소비하는지 검사한다.
  \item Singleton/overlap/boundary/hard-attachment/traveling-support에서
  support-restricted per-RHS route/joint Schur, adjustment/reaction amplification,
  generalized work, halo cap과 free/hard/compliant post-step \(r_{\mathrm{ROM}}\) ledger를 검증한다.
\end{enumerate}

Canonical source label은 각각
\code{eq:relation-mls-map}, \code{eq:canonical-structural-force-ledger},
\code{eq:klin-projective-hard-gate}, \code{eq:v0-structural-freeze-signature},
\code{eq:fixed-global-sample-jacobian}, \code{eq:support-restricted-global-routing},
\code{eq:hard-reaction-rom-residual-decomposition}다.

\subsection{Critical path}

\begin{lstlisting}
TD00 governance
  -> TD01 deterministic contracts/E0 + Step 0 preflight
      |-> minimal independent-GS schema/relation probe [parallel fixture only]
      `-> after Step-0 gate pass
          |-> TD02 teacher + independent GS data
          |   -> TD03 oracle object package
          |       -> TD04 full-anchor current-surface aero
          |       -> TD05 Global ROM + same-anchor kill test [MVP-A]
          |       -> TD07-A patch/local operator + single-patch unit solver
          |       -> TD11-A unique-anchor assembly + joint complement
          |       -> TD07-B Oracle coupled Local + Gate B [MVP-B]
          |       -> TD10-core cached-compliance selector/state machine [MVP-C]
          |       -> TD11-B corrected-aero/structural-delta feedback
          |       -> TD08 projection/labels
          |           |-> TD09 optional F missing-force expert
          |           `-> TD10-G optional learned gate
          `-> TD12-R formal 5--10-object representation feasibility [parallel]
              -> TD12 full predicted scaffold/package after TD12-R and
                 Oracle downstream-sensitivity checks [MVP-D]
  -> TD13 target runtime + GS transport
  -> TD14 comparison experiments/paper evidence

Optional branch after TD05 profiling:
  TD06 optional H aerodynamic hyper-reduction
\end{lstlisting}

여기서 Step-0 gate는 K00--K09의 deterministic 해당 항목, K20, K22--K23과
K21의 schema/synthetic-aggregator/fallback fixture를 뜻한다. Learned K21 calibration은
Track B0/E1에서 별도로 판정하며, K10--K19와 K24는 각 downstream stage에서 판정한다.
Step 0 중에는 fixture에 필요한 최소 Oracle strip/flag와 값싼 independent-GS
schema/relation probe만 병렬 허용한다. Formal 5--10-object TD12-R,
full teacher/data 생성과 Track A0는 Step-0 gate 통과 뒤 병렬 진행한다.
TD06은 TD07 진입 조건이 아니며 full-anchor backend가 V0 reference이자 fallback이다.
TD12의 full topology network는 TD12-R과 TD05/TD07/TD10-core/TD11의 Oracle-package downstream sensitivity가 확인된 뒤 본격 학습한다.
공식 Gate B 결과는 TD11-A의 joint conservation--complement를 거친 force와 single coupled Local solve로만 만든다.
TD11-B one-corrector는 Gate B 뒤의 SYS ablation이며 Local necessity를 gate하지 않는다.

\subsection{Milestone overview}

\small
\begin{longtable}{L{0.065\linewidth}L{0.235\linewidth}L{0.47\linewidth}L{0.075\linewidth}}
\toprule
ID & Milestone & 핵심 완료 기준 & 상태 \\
\midrule
\endfirsthead
\toprule
ID & Milestone & 핵심 완료 기준 & 상태 \\
\midrule
\endhead
TD00 & Project governance/reset & 새 method용 config, run manifest, artifact 경로와 상태 규칙이 재현된다. & \done \\
TD01 & Deterministic contracts/E0 & Typed schema, force ledger, canonical-reference manifest와 basis/operator/transport P0 unit tests가 통과한다. & \todo \\
TD02 & Teacher/data pipeline & Force breakdown, split, independent GS가 포함된 누수 없는 dataset이 생성된다. & \todo \\
TD03 & Oracle object package & Oracle anchors/operators/bases/patches/binding package가 validation을 통과한다. & \todo \\
TD04 & Current-surface full aero & Module 1--3 full-anchor evaluator가 rigid/relative-wind/normal-area tests를 통과한다. & \todo \\
TD05 & Global ROM predictor & Start/stop/chirp에서 frequency, damping, energy와 same-anchor baseline을 검증한다. & \todo \\
TD06 & Optional H aero reduction & Full-anchor가 병목일 때만 generalized-load fidelity와 wall-clock 이득을 검증한다. & \todo \\
TD07 & Complementary Local physics & TD11-A joint force를 사용하는 Oracle active set에서 coupled Local solve가 same-p95 Global rank sweep보다 stable high-frequency 이득을 보인다. & \todo \\
TD08 & Projection/label generation & Generalized-to-nodal route, positive-benefit와 feedback label의 차원/owner consistency가 확인된다. & \todo \\
TD09 & Optional F missing-force expert & Always-on F가 physics-only Local보다 rollout/spectrum을 개선할 때만 채택한다. & \todo \\
TD10 & Fixed-\(K\) selector / optional G & Cached-compliance/Oracle fixed-cardinality selector를 완성하고, G는 같은 grouping/\(K_A\)와 measured-latency Pareto를 넘을 때만 채택한다. & \todo \\
TD11 & Assembly/complement/feedback & A: unique-anchor small-Schur joint assembly를 Gate B 전에 고정하고, B: Gate B 뒤 one-corrector SYS test를 통과한다. & \todo \\
TD12 & Predicted target package & Static GS와 setup metadata에서 explicit target triangle-mesh reconstruction 없이 package를 만들고 Oracle gap을 제한한다. & \todo \\
TD13 & Integrated runtime/transport & 정확한 14-module trace, affine GS transport, p50/p95 profile과 rendering이 재현된다. & \todo \\
TD14 & Evaluation/paper evidence & E0--E10, baselines, OOD, statistics와 go/no-go 결과가 완결된다. & \todo \\
\bottomrule
\end{longtable}
\normalsize

\subsection{Runtime Module과 milestone 대응}

\begin{longtable}{L{0.05\linewidth}L{0.27\linewidth}L{0.21\linewidth}L{0.29\linewidth}}
\toprule
M & 기능 & 주 milestone & 다음 typed 출력 \\
\midrule
\endfirsthead
\toprule
M & 기능 & 주 milestone & 다음 typed 출력 \\
\midrule
\endhead
1 & 이전 state/version 읽기 & TD01, TD13 & \code{PreparedFrameState} \\
2 & Current anchor/surface 복원 & TD04, TD13 & \code{SurfaceState} \\
3 & Relative wind/base aero & TD04 & \code{AnchorAeroForces} \\
4 & Full generalized-load projection / optional H & TD04, TD06 & \code{ReducedAeroLoad} \\
5 & Always-on Global predictor & TD05 & \code{GlobalPrediction} \\
6 & Fixed-\(K\) selector / optional G gate & TD07, TD10 & \code{SelectorDecision} \\
7 & Active Local base + optional F proposal & TD04, TD07, TD09 & \code{PatchForceBatch} \\
10 & Overlap-aware assembly & TD11 & \code{AssemblyResult} \\
9 & Global force span 제거 & TD07, TD11 & \code{ComplementForce} \\
8 & Single Local physical solve & TD07, TD11 & \code{LocalPrediction} \\
11 & Local-to-Global corrector & TD11 & \code{CorrectedDynamicState} \\
12 & Final anchor state & TD13 & \code{FinalAnchorState} \\
13 & Gaussian transport/render & TD13 & \code{DeformedGaussianAsset} \\
14 & State machine/diagnostics & TD10, TD13 & \code{RuntimeState}, \code{FrameDiagnostics} \\
\bottomrule
\end{longtable}

\section{공통 artifact registry}

모든 milestone은 code와 함께 아래 산출물을 남긴다.

\begin{longtable}{L{0.075\linewidth}L{0.31\linewidth}L{0.52\linewidth}}
\toprule
ID & Artifact & 최소 내용 \\
\midrule
\endfirsthead
\toprule
ID & Artifact & 최소 내용 \\
\midrule
\endhead
A00 & \artifact{run_manifest.json} & Git commit, config hash, seed, device, package/dataset version, command. \\
A01 & \artifact{contract_report.json}, \artifact{v0_contract_manifest.json} & Units/frames/shapes, force ledger, PSD/symmetry/orthogonality/transport tests와 canonical source/hash/label/API/owner reference validation. \\
A02 & \artifact{teacher_manifest.json} & Asset/material/attachment/wind/trajectory/split/force-channel inventory. \\
A03 & \artifact{oracle_object_package/} & Anchors, graph/local relations, patch-indexed bases/cross blocks, optional H placeholder, binding, validation. \\
A04 & \artifact{full_aero_report.json} & Relative-wind, normal/area, force component, net wrench unit tests. \\
A05 & \artifact{global_rollout_report.md} & Frequency, amplitude, phase, damping, energy, latency와 failure cases. \\
A06 & \artifact{aero_reduction_report.md} & Full-anchor profile과 조건부 H/random/FPS/RSH/direct accuracy--latency; generalized load를 primary로 보고. \\
A07 & \artifact{local_physics_report.md} & Basis complement, state transfer, energy, oracle-active Local Pareto. \\
A08 & \artifact{training_labels/} & Projected state/feedback delta와 조건부 F missing-force, 조건부 G positive-benefit label, confidence. \\
A09 & \artifact{optional_F_missing_force_model/} & F를 채택한 경우의 config, checkpoint, normalization, validation/rollout report. \\
A10 & \artifact{selector_report/} & Fixed-\(K\) Oracle/cached-compliance selector, K sweep, full-counter analysis, optional G checkpoint/calibration/cost extension과 dispatcher profile. \\
A11 & \artifact{coupling_report.md} & Dense--Schur equivalence, constraint rank/condition, factor share/rebuild와 latency share, overlap seam, force/torque, leakage, corrector double-count와 stability. \\
A12 & \artifact{predicted_object_package/} & TD12-R early representation report, Static-GS inferred package, confidence/fallback, Oracle/kNN gap. \\
A13 & \artifact{runtime_report.md} & Exact trace, p50/p95, memory, kernel count, renders, invalid-state counts. \\
A14 & \artifact{paper_evidence/} & E0--E10 tables, plots, videos, statistics, baseline manifests. \\
\bottomrule
\end{longtable}

\newpage
\section{TD00. Project governance와 method reset}

\textbf{목표.} 이전 proxy/SH-residual 방향과 새 Global--Local physics 방향의 상태가 섞이지 않도록
새 experiment namespace, source-of-truth, run manifest와 변경 기록 규칙을 고정한다.

\textbf{선행 조건.} 없음.

\textbf{입력.}
현재 아이디어 스케치, 이 체크리스트, 기존 \path{code/}, \path{experiments/} 자산.

\subsection*{구현 체크}

\begin{itemize}
  \item[\done] 이전 문서와 기존 implementation checklist를 backup에 원래 이름으로 보존한다.
  \item[\done] 현재 아이디어 source/PDF/bibliography/bundle을 \path{ideas/} 최상위에 둔다.
  \item[\done] 새 체크리스트 파일명을 현재 method title에 맞게 고정한다.
  \item[\done] \path{experiments/TD00_contracts/README.md}를 만든다.
  \item[\done] 공통 experiment README template에 다음 필드를 넣는다.
    \begin{itemize}
      \item goal과 falsification question,
      \item exact input privilege와 runtime restrictions,
      \item command/config/seed/device,
      \item outputs/metrics/completion gate,
      \item failed cases와 decision.
    \end{itemize}
  \item[\done] 모든 run이 \artifact{run_manifest.json}을 기록하도록 공통 utility를 만든다.
  \item[\done] config가 dataset/package/model hash를 저장하도록 한다.
  \item[\done] code-side, idea-side, experiment-side session 기록 책임을 README에 명시한다.
  \item[\done] 기존 \code{M01--M03} 구현을 새 완료 상태로 승격하는 audit checklist를 만든다.
  \item[\done] RSH/appearance SH 코드를 main runtime dependency로 import하지 않는지 dependency scan을 추가한다.
\end{itemize}

\subsection*{필수 산출물}

\begin{itemize}
  \item \artifact{experiments/TD00_contracts/README.md}
  \item \artifact{run_manifest.schema.json} 또는 동등한 typed schema
  \item Reuse-candidate audit report
  \item 새 milestone tag를 사용하는 최소 smoke-test run
\end{itemize}

\textbf{완료 기준.}
빈 smoke run 하나가 source commit, config, seed, package/dataset version, device와 output 경로를 자동 기록하고,
이전 method code가 새 completion status에 섞이지 않는다.

\textbf{완료 증거.}
Python 3.12 unit test 38개와 clean-source reference run
\code{td00-contracts-smoke-20260812t184509388327z-abde7ba}가 통과했다.
Compact manifest/report/completion marker의 hash를 독립 verifier로 재검증했다.
Python 3.10/3.11 실제 실행과 물리 E0는 TD01 이후 미검증 항목으로 유지한다.

\section{TD01. Deterministic schema, 좌표계, force ledger와 E0}

\textbf{목표.}
학습을 시작하기 전에 모든 물리량의 shape, dtype, device, 단위, 좌표계, index, owner와 frame dependency를 코드로 고정한다.

\textbf{선행 조건.} TD00 run/artifact convention.

\textbf{입력.}
Canonical GS fixture, 작은 oracle strip fixture, 현재 method의 state/force ownership 정의.

\subsection*{필수 typed data}

\begin{lstlisting}
CanonicalGaussianAsset
TeacherTrajectory
AnchorScaffold
OperatorRelationSetQop
RelationValidityReport / AnchorAbsoluteConfidence
RestLinearOperators             # M, K_lin^0, C_lin^0; modal/calibration reference
ProjectiveRuntimeOperators       # K^0=K_proj^0, C^0=C_proj^0; canonical runtime
ProjectiveTargetLedger           # frozen targets/rates, h_K, h_C, fit/timing lineage
ReducedOperators
PatchBasisBlock / PatchCrossBlocks
SurfaceTangentMap3x2
TangentAffineMapRank2_3x3
FullShellMap3x3
AeroCubature                    # optional H only
GaussianBinding
ObjectPackage
RuntimeState
PreparedFrameState
SurfaceState
AnchorAeroForces
ReducedAeroLoad
FixedGlobalSampleLifting         # setup-fixed T_s and Phi_s=T_s Phi
GlobalPrediction
SelectorDecision
PatchForceBatch
ConstraintOperator
ConstraintSignature
ConstraintSolveReport
SupportRoutingReport
AssemblyResult
ComplementForce
LocalPrediction
CorrectedDynamicState
SupportReactionLedger            # free/hard/compliant + post-step r_ROM
FinalAnchorState
DeformedGaussianAsset
FrameDiagnostics
FrameOutput
CanonicalContractReference       # source path/hash, labels, API/owner/schema versions
\end{lstlisting}

\subsection*{Schema와 convention 체크}

\begin{itemize}
  \item[\todo] World, object, material anchor, local patch, camera frame의 transform 방향을 문서화한다.
  \item[\todo] Right-handed convention, front/back, outward normal과 two-sided sheet sign을 고정한다.
  \item[\todo] Length/time/mass/force 단위를 SI로 둘지 nondimensional preset으로 둘지 artifact마다 명시한다.
  \item[\todo] Metric-scale asset과 normalized asset의 혼용을 validator가 거부하도록 한다.
  \item[\todo] Tensor layout, dtype, device, batch/anchor/patch/sample index convention을 고정한다.
  \item[\todo] Quaternion order와 covariance parameterization을 고정한다.
  \item[\todo] Package, state, model, dataset version과 backward compatibility 규칙을 정의한다.
  \item[\todo] TD01이 \artifact{v0_contract_manifest.json}을 생성한다. 여기에는
  \code{canonical\_source\_path}, clean commit 또는 dirty-content hash,
  \code{canonical\_labels[]}, typed API/owner/schema version, 동결한 V0 branch와 failure policy를 저장한다.
  Canonical TeX를 parse해 각 required label이 정확히 한 번 존재하는지 검사하고,
  stale label/API/owner/version/hash가 있으면 fail-fast 한다. TD00의 완료 상태는 다시 열지 않는다.
  \item[\todo] Setup metadata와 frame runtime input을 별도 schema로 두고, frame 중 material/attachment layout 변경을 validator가 거부하도록 한다.
  \item[\todo] Patch ID, fixed coordinate range, overlap group과 owner weight를 package에 고정한다.
  \item[\todo] \(\mathcal Q_{\mathrm{op}}\) relation set, channel별 calibrated validity,
  anchor absolute confidence와 support-expand/fallback/reject reason을 서로 다른 typed field로 둔다.
  Area partition logit, reconstruction confidence, relation validity와 material coefficient를 alias하지 않는다.
  \item[\todo] Rest-linear \((K_{\mathrm{lin}}^0,C_{\mathrm{lin}}^0)\)와 canonical runtime
  \((K^0,C^0)=(K_{\mathrm{proj}}^0,C_{\mathrm{proj}}^0)\), frozen target ledger를 별도 type/version으로 둔다.
  \item[\todo] Setup-fixed sample extraction \(T_s\)와
  \(\Phi_s:=T_s\Phi\)를 package에 저장하고 frame state가 Jacobian/basis를 바꾸지 못하게 한다.
  \item[\todo] Surface tangent 3\(\times\)2 map과 covariance용 full 3\(\times\)3 map을 서로 다른 typed field/API로 둔다.
  \item[\todo] Save--load round trip 후 값, shape, dtype, index, units가 동일한지 검사한다.
  \item[\todo] State와 package version/hash mismatch는 즉시 fail-fast 한다.
  \item[\todo] \code{active\_set}, \code{decay\_set}, \code{dormant\_set}, diagnostics index가 서로 섞이지 않게 한다.
\end{itemize}

\subsection*{Force ledger 체크}

\begin{itemize}
  \item[\todo] 한 frame에 들어오는 모든 force channel의 owner, frame, sample/anchor/generalized space를 기록한다.
  \item[\todo] Gravity, attachment, base aero, structural base, missing force, feedback delta를 enum/channel로 구분한다.
  \item[\todo] Constraint로 소비된 force와 다음 corrector가 소비할 delta를 기록한다.
  \item[\todo] 같은 channel이 predictor와 corrector에서 두 번 더해지면 test가 실패하도록 한다.
  \item[\todo] Teacher residual label 생성 시 제거한 base channel 목록을 label metadata에 저장한다.
  \item[\todo] Structural restoring은 generic \(-Ku-Cv\) 한 channel로 축약하지 않는다.
  Canonical runtime에서 \(-K^0u+h_K\)와 \(-C^0v+h_C\)를 각각 한 번 소비하고,
  target/target-rate, fit state, hold 시점과 source label을 \code{ProjectiveTargetLedger}에 저장한다.
  \item[\todo] \(K_{\mathrm{lin}}^0,C_{\mathrm{lin}}^0\)은 modal/calibration 및 rest-linear ablation owner이고,
  canonical rollout force ledger에 \(K^0,C^0\)와 동시에 더해지지 않게 한다.
  \item[\todo] External support-routed Global component, Local remainder, joint-KKT adjustment,
  free/hard companion, compliant object force/handle reaction과 post-step total hard reaction을 서로 다른 channel로 둔다.
  \item[\todo] Post-step residual은 implicit-midpoint timing의 final state에서만 계산하고,
  hard reaction이 reduced truncation residual \(r_{\mathrm{ROM}}\)을 흡수하지 않게 별도 저장한다.
\end{itemize}

\subsection*{E0 deterministic unit tests}

\begin{itemize}
  \item[\todo] \(M\) positive definite와
  \(K_{\mathrm{lin}}^0,C_{\mathrm{lin}}^0,K^0,C^0\)의 typed symmetry/PSD test.
  Rest-linear와 projective runtime operator를 교환하거나 이중 적용하면 실패하는 negative test.
  \item[\todo] Reduced block dimension, symmetry와 reciprocity test.
  \item[\todo] \(\Phi^\top M\Phi\approx I\), patch별 \(\Phi^\top M\Psi_r\approx0\), patch owner 보존 test.
  \item[\todo] Reachable patch/group의 \(M_{rs},C_{rs},K_{rs}\) block reciprocity, rank와 condition-number test.
  \item[\todo] Zero force/zero wind에서 rest가 invariant인지 test.
  \item[\todo] Rigid translation/rotation에서 internal elastic energy가 불필요하게 생기지 않는지 test.
  \item[\todo] Pure-rigid 및 strained-state superposed-rotation에서 projective energy/force equivariance,
  relative-rate damping, planar in-plane bending pollution과 one-shot/two-step/converged target-splitting test.
  \item[\todo] Held-target stiffness/damping action, action-tangent symmetry/PSD,
  first 8--16 mode/subspace와 common-probe closure test.
  \item[\todo] Canonical implicit midpoint 대 backward Euler/rest-linear ablation의 phase,
  target-band PSD, damping envelope, energy와 timestep convergence test.
  \item[\todo] Constant force와 impulse에 대한 analytic 작은-system response test.
  \item[\todo] Identity/translation/rotation/affine Gaussian reproduction과 covariance SPD test.
  \item[\todo] Rest-chart normalization, rigid rotation, in-plane scale/shear에서 3\(\times\)2 cross-product area-normal과 \(F_\parallel=J^t(J^0)^\dagger\) test.
  \item[\todo] 3\(\times\)2 tangent map을 3\(\times\)3 covariance transform에 잘못 전달하면 실패하는 type/negative test.
  \item[\todo] Net force/torque 계산의 center convention test.
  \item[\todo] Current-geometry torque operator와 center translation identity test.
  \item[\todo] Small fixture에서 dense primal KKT와 unique-anchor small-Schur의 force/objective/constraint residual equivalence test.
  \item[\todo] Duplicate/dependent constraint row의 consistent/inconsistent RHS rank test.
  \item[\todo] 같은 support operator에서 external/structural channel의 RHS를 각각 검사하고,
  aggregate feasible rate만으로 infeasible RHS를 숨기지 않는 per-RHS consistency test.
  \item[\todo] Singleton/overlap/boundary/hard-attachment/traveling-support에서
  support expansion과 halo cap, raw-force adjustment, reaction amplification,
  generalized-work 보존 및 post-step \(r_{\mathrm{ROM}}\) 분해 test.
  \item[\todo] Geometry/support/row signature가 바뀐 factor를 재사용하거나 sequential projection으로 fallback하면 실패하는 negative trace test.
  \item[\todo] CPU/GPU 또는 float32/float64 reference 차이를 기록한다.
  \item[\todo] 고정 seed에서 같은 trace와 artifact hash가 나오는지 test.
  \item[\todo] Canonical label 삭제/중복, API owner/version 또는 source hash mismatch를 주입하면
  derived-contract reference test가 정확히 실패하는지 검사한다.
\end{itemize}

\subsection*{권장 API}

\begin{lstlisting}[language=Python]
validate_object_package(package) -> ContractReport
validate_runtime_state(state, package, handle_motion, dt) -> PreparedFrameState
validate_force_ledger(ledger) -> ForceOwnershipReport
run_e0_contract_suite(fixtures, tolerances) -> ContractReport
validate_canonical_contract_reference(manifest, canonical_tex) -> ContractReport
\end{lstlisting}

\subsection*{필수 산출물}

\begin{itemize}
  \item \artifact{contract_report.json}
  \item \artifact{v0_contract_manifest.json}
  \item Canonical label/API/owner/schema stale-reference negative-test report
\end{itemize}

\textbf{완료 기준.}
\artifact{contract_report.json}에 모든 Step-0-relevant P0 deterministic check가 통과하고,
\artifact{v0_contract_manifest.json}의 canonical reference validation이 통과하며,
schema/version/owner 오류를 의도적으로 주입한 negative test가 정확히 실패한다.
이 milestone 전에는 teacher label이나 learned model을 production dataset으로 만들지 않는다.

\section{TD02. Teacher, independent GS와 dataset split}

\textbf{목표.}
Global/Local solver와 topology distillation을 검증하고,
채택 시 조건부 F/G도 학습할 수 있는 누수 없는 teacher dataset을 만든다.

\textbf{선행 조건.} TD01 data schema, units/frames와 force ledger.

\textbf{입력.}
Source thin-shell mesh, explicit material metadata/attachment layout, prescribed wind program, render cameras와 reconstruction config.

\subsection*{Teacher-S structural trajectory}

\begin{itemize}
  \item[\todo] 첫 fixture를 한쪽 고정 cantilever strip/flag로 고정한다.
  \item[\todo] 후속 asset으로 two-point cloth, curtain strip, paper sheet, leaf를 준비한다.
  \item[\todo] High-resolution nonlinear shell/FEM 또는 검증된 dense MPM backend를 선택한다.
  \item[\todo] Timestep, substep, integrator, material law와 damping convention을 manifest에 기록한다.
  \item[\todo] Position, velocity, acceleration을 저장한다.
  \item[\todo] Strain, curvature, stretch/bending energy를 저장한다.
  \item[\todo] Surface normal, area, deformation gradient를 저장한다.
  \item[\todo] Internal, aerodynamic, gravity, attachment/reaction force breakdown을 가능한 한 solver에서 직접 저장한다.
  \item[\todo] 저장된 force breakdown으로 teacher acceleration을 다시 구성하는 balance test를 추가한다.
  \item[\todo] Optional Gaussian mean/covariance target과 multi-view render를 저장한다.
\end{itemize}

\subsection*{Wind와 setup-metadata sweep}

\begin{itemize}
  \item[\todo] Wind direction/elevation과 weak--strong magnitude sweep.
  \item[\todo] Steady, sinusoidal gust, chirp, abrupt start/stop.
  \item[\todo] Stationary, slow-traveling, fast-traveling, reversing와 localized/nonuniform spatial gust.
  \item[\todo] Setup에 사용자가 지정하는 density, thickness, stretch/shear/bending stiffness, damping, anisotropy.
  \item[\todo] Setup에 사용자가 지정하는 one edge, two points, partial/compliant support, stem attachment layout.
  \item[\todo] Attachment layout 또는 basis를 바꾸는 material 값이 달라지면 새 package/version으로 분리한다.
  \item[\todo] Frame 중에는 고정 support DOF의 prescribed handle motion만 바꾸고 layout edit는 수행하지 않는다.
  \item[\todo] Rest, pre-deformed, initial-velocity condition.
  \item[\todo] 각 factor와 random seed가 trajectory manifest에서 검색 가능하도록 한다.
\end{itemize}

\subsection*{Optional Teacher-FSI track}

\begin{itemize}
  \item[\todo] CFD/LBM 또는 계측 subset의 surface sample/force schema를 Teacher-S와 정렬한다.
  \item[\todo] 같은 geometry에서 quasi-steady와 FSI force를 저장한다.
  \item[\todo] \code{unsteady\_defect = FSI/measurement - quasi\_steady}를 명시적으로 만든다.
  \item[\todo] Teacher-FSI가 없는 sample에는 wake/vortex/history label을 생성하지 않는다.
  \item[\todo] Teacher-S-only와 Teacher-FSI sample을 dataset manifest에서 분리한다.
\end{itemize}

\begin{tcolorbox}[colback=LightOrange,colframe=orange!60!black,title=Claim 제한,fonttitle=\bfseries]
Teacher-FSI 또는 실측 force가 없으면 learned force는 basis truncation, nonlinear structure,
operator/topology distillation defect만 보정한다고 기술한다. Wake, vortex shedding 또는 unsteady aerodynamics를 학습했다고 주장하지 않는다.
\end{tcolorbox}

\subsection*{Independent static-GS variants}

\begin{itemize}
  \item[\todo] Source mesh에 Gaussian을 직접 붙인 oracle/debug asset을 별도 표시한다.
  \item[\todo] Rendered views로부터 target static GS를 독립적으로 reconstruct한다.
  \item[\todo] View count, reconstruction seed, Gaussian density variant를 만든다.
  \item[\todo] Hole, floater, covariance noise, opacity perturbation, missing backside variant를 만든다.
  \item[\todo] Teacher와 static GS의 coordinate/metric scale alignment를 저장한다.
  \item[\todo] Alignment residual이 tolerance를 넘는 asset은 fail/fallback으로 표시한다.
\end{itemize}

\subsection*{Split와 leakage 방지}

\begin{itemize}
  \item[\todo] Frame-random split을 금지하고 trajectory 단위 split을 구현한다.
  \item[\todo] Object-disjoint split을 별도로 만든다.
  \item[\todo] Relation-validity와 anchor absolute-confidence calibrator를 fit/freeze할
  development-calibration split도 final test object와 object-disjoint하게 만든다.
  \item[\todo] Calibration/coverage report를 anchor density, candidate count, valence와
  MLS condition-number bin으로 층화한다. 한 density/valence의 aggregate ECE만 보고 통과시키지 않는다.
  \item[\todo] Unseen wind/setup-material/setup-attachment/geometry/category/reconstruction split을 만든다.
  \item[\todo] Combined object+material+wind OOD split을 만든다.
  \item[\todo] 같은 source trajectory의 GS variants가 train/test에 흩어져 identity leakage를 만들지 않게 group split한다.
\end{itemize}

\subsection*{필수 산출물}

\begin{lstlisting}
datasets/
  manifests/teacher_manifest.json
  teacher_structure/<trajectory_id>/
  teacher_fsi/<trajectory_id>/          # optional
  canonical_gs/<asset_id>/
  gs_variants/<asset_id>/<variant_id>/
  alignment/<asset_id>/<variant_id>.json
  splits/{train,val,calibration,test,ood_*.json}
\end{lstlisting}

\textbf{완료 기준.}
최소 flag/strip의 start/stop, sinusoidal, chirp, traveling gust trajectory가 force breakdown과 함께 재생 가능하고,
force balance, coordinate alignment, trajectory/object split leakage test가 통과한다.

\section{TD03. Oracle runtime object package와 reduced bases}

\textbf{목표.}
Predicted topology의 오차를 배제한 clean oracle case에서 runtime data structure와 Global/Local physical space를 먼저 닫는다.

\textbf{선행 조건.} TD01 contracts, TD02 oracle source/trajectory.

\textbf{입력.}
Source mesh topology, user material metadata/attachment layout, canonical/independent GS, teacher snapshots.

\subsection*{Oracle AnchorScaffold}

\begin{itemize}
  \item[\todo] Material-space anchor 위치, rest normal, area, mass와 reliability를 만든다.
  \item[\todo] Oriented material edges와 bending/hinge relations를 저장한다.
  \item[\todo] Hard/compliant attachment를 하나의 명시적 representation으로 저장한다.
  \item[\todo] Core--halo patch와 overlap graph를 만든다.
  \item[\todo] Partition-of-unity로 total area와 mass가 원본과 일치하는지 검사한다.
  \item[\todo] Disconnected close layer를 일부러 포함한 negative fixture를 만든다.
\end{itemize}

\subsection*{Physical operators}

\begin{itemize}
  \item[\todo] \code{OperatorRelationSetQop}에 incidence support, rest frame,
  metric MLS stencil, membrane/bending area share, channel별 validity와 condition metadata를 저장한다.
  Canonical source는 \code{eq:relation-mls-map},
  \code{eq:relation-area-partition}과 \code{eq:relation-absolute-confidence}다.
  \item[\todo] Positive diagonal/lumped \(M\)과 modal/calibration reference인
  \(K_{\mathrm{lin}}^0,C_{\mathrm{lin}}^0\)를 별도 조립한다.
  \(K_{\mathrm{lin}}^0\)은 stretch/shear/bending의 rest-linear
  \(B^\top D B\), \(C_{\mathrm{lin}}^0\)은 channel별 Kelvin--Voigt reference이며,
  generic Rayleigh damping을 canonical internal force에 넣지 않는다.
  \item[\todo] Runtime에는 world-vector projective relation과 objective normal-curvature target으로
  \(K^0:=K_{\mathrm{proj}}^0\), \(C^0:=C_{\mathrm{proj}}^0\)를 조립한다.
  Canonical source는 \code{eq:objective-vector-projective-constraints},
  \code{eq:bending-normal-projectors}, \code{eq:canonical-structural-force-ledger}다.
  \item[\todo] Frame별 frozen rotation/normal-curvature target과 target-rate에서
  \(h_K,h_C\)를 만들고 \code{ProjectiveTargetLedger}에 fit state, hold time과
  operator/package version을 저장한다. Runtime restoring은
  \(-K^0u+h_K-C^0v+h_C\)이며 target RHS를 두 번 더하지 않는다.
  \item[\todo] Hard attachment는 free-DOF restriction/lift로, compliant attachment는
  명시적 object-force/handle-reaction channel로 처리한다. 같은 support를 stiffness, projection과
  external attachment force에 중복 소유시키지 않는다.
  \item[\todo] \(M\), rest-linear와 projective operator 각각의 PSD/symmetry,
  attachment null-space, rigid/equivariance, normal-curvature pollution과 relative-rate damping test를 실행한다.
  \item[\todo] Held action과 locally refit energy/force/action tangent, modal subspace와
  common probe closure를 검사한다. Action-validity와 rest-linear-agreement를 별도 결과로 저장하되,
  canonical package는 둘을 모두 통과해야 한다. PSD/objectivity/action 자체가 실패하면 backend를 재설계하고,
  modal/probe agreement만 실패한 symmetric action-tangent basis는 기존 gate의 우회 통과가 아니라
  MC1 claim과 모든 basis/complement/cache를 바꾸는 별도 versioned pivot으로만 검토한다.
  One-shot splitting만 실패하면 two-step/converged target fit을 진단 대안으로 검토하고,
  그래도 실패할 때만 state-dependent co-rotated tangent와 factor/latency 계약을 새로 검증한다.
  \item[\todo] Setup material/attachment 변경 시 어떤 block, basis와 factorization을 invalidate하고 새 package version을 만드는지 기록한다.
\end{itemize}

\subsection*{Global/Local bases}

\begin{itemize}
  \item[\todo] Canonical V0는 rest-linear \(K_{\mathrm{lin}}^0\) eigenmode와
  physics/POD hybrid 후보에서 \(\Phi\)를 만든다. Rest-linear agreement가 실패한
  action-tangent eigenbasis는 canonical closure를 통과한 것으로 부르지 않고,
  MC1 operator/eigenspace claim을 철회한 별도 versioned pivot에서만 허용한다.
  Raw \(K_{\mathrm{proj}}^0\)를 thin-shell tangent로 간주해 basis를 자동 생성하지 않는다.
  \item[\todo] \(\Phi^\top M\Phi=I\)를 검사한다.
  \item[\todo] Patch raw basis \(B_r\)와 patch-to-global scatter를 만든다.
  \item[\todo] \(U_r=R_r^\top W_rB_r\)로 anchor space에 lift하되 patch ID와 column owner를 유지한다.
  \item[\todo] 각 patch에서만 \(\bar B_r=P_\perp U_r\),
  \(P_\perp=I-\Phi(\Phi^\top M\Phi)^{-1}\Phi^\top M\)로 Global span을 제거한다.
  \item[\todo] Patch Gram \(G_r=\bar B_r^\top M\bar B_r=V_r\Lambda_rV_r^\top\)를 고유분해하고,
  \(\lambda_i>\tau_{\mathrm{rank}}\lambda_{\max}\)인 retained eigenspace \((V_{r,+},\Lambda_{r,+})\)에서만
  \(\Psi_r=\bar B_rV_{r,+}\Lambda_{r,+}^{-1/2}\)로 정규화한다.
  제거된 near-null column, threshold와 retained rank를 patch metadata에 저장한다.
  \item[\todo] 모든 patch를 한 번에 QR/SVD하여 서로 다른 patch basis를 섞는 code path를 negative test로 금지한다.
  \item[\todo] Patch별 \(\Phi^\top M\Psi_r=0\), \(\Psi_r^\top M\Psi_r=I\) tolerance와 fixed coordinate range를 저장한다.
  \item[\todo] 서로 다른 patch끼리는 직교화를 강제하지 않고
  \(M_{rs}=\Psi_r^\top M\Psi_s\), \(C_{rs}=\Psi_r^\top C\Psi_s\),
  \(K_{rs}=\Psi_r^\top K\Psi_s\)를 저장한다.
  \item[\todo] Strong overlap로 superset Gram이 rank deficient하면 overlap-connected patch를 하나의 gating group으로 묶거나 owner-preserving block pivoting으로 중복 column을 제거한다.
  \item[\todo] V0는 \(\Psi_{\mathrm{sup}}=[\Psi_1,\ldots,\Psi_P]\)의 fixed slot과 persistent full state
  \(z_{\mathrm{sup}},\dot z_{\mathrm{sup}}\)를 사용한다.
  \(S_t=A_t\cup D_t\) selector matrix \(E_{S_t}\)로
  \(z_{S_t}=E_{S_t}^\top z_{\mathrm{sup}}\)를 gather하고 active principal block을 한 번 푼 뒤,
  같은 fixed slot에 scatter한다. Frame마다 basis dimension을 바꾸거나 state를 remap하지 않는다.
  \item[\todo] Reachable active/decay group마다 reduced block rank/conditioning, symmetry와 reciprocity를 검사한다.
  \item[\todo] \(K_{gr},C_{gr}\)가 수치적으로 0인지 nonzero인지 기록하고 corrector claim과 연결한다.
\end{itemize}

\begin{tcolorbox}[colback=LightOrange,colframe=orange!60!black,title=Patch energy와 계산 단위,fonttitle=\bfseries]
Inter-patch cross block 때문에 전체 Local energy는 patch별 scalar의 단순 합이 아니다.
State-machine 종료에는 owner-weighted support energy 또는 overlap-group energy를 사용하고,
정의와 cross-term 배분을 package에 저장한다. Patch 수, 선택 단위, coordinate 수와 실제 solver cost의 관계를 함께 profile한다.
Dynamic basis recomputation/state remapping은 V0 범위가 아니다.
\end{tcolorbox}

\subsection*{Gaussian binding}

\begin{itemize}
  \item[\todo] 각 Gaussian을 고정된 소수 anchor support와 partition weight에 bind한다.
  \item[\todo] Rest offset와 optional local frame을 저장한다.
  \item[\todo] Frame마다 nearest anchor를 다시 고르지 않게 한다.
  \item[\todo] Identity/translation/rotation/affine reproduction을 TD01 suite와 연결한다.
\end{itemize}

\subsection*{권장 API와 package}

\begin{lstlisting}[language=Python]
build_oracle_object_package(
    source_surface, canonical_gs, material, attachment, config
) -> ObjectPackage

save_object_package(package, path)
load_object_package(path) -> ObjectPackage
validate_object_package(package) -> ContractReport
\end{lstlisting}

\begin{lstlisting}
oracle_object_package/
  package_manifest.json
  anchors.npz
  topology.npz
  operators.npz
  bases.npz
  patches.npz
  aero_samples_placeholder.npz
  gaussian_binding.npz
  validation.json
\end{lstlisting}

\textbf{완료 기준.}
Oracle strip package가 save--load 후 동일하며, operator PSD/symmetry, basis orthogonality,
area/mass conservation, binding reproduction test를 모두 통과한다.

\section{TD04. Current-surface state와 full aerodynamic reference}

\textbf{목표.}
Runtime Module 1--3을 learning/hyper-reduction 없이 구현하여 current geometry와 relative wind에서 full-anchor aerodynamic force를 얻는다.

\textbf{선행 조건.} TD03 oracle package, TD01 transport/force convention.

\subsection*{Module 1: state preparation}

\begin{itemize}
  \item[\todo] \(q,\dot q,z,\dot z\), active/decay state, memory와 package version을 읽는다.
  \item[\todo] Dimension/index/device/version mismatch를 fail-fast 한다.
  \item[\todo] Setup에서 고정된 support DOF의 prescribed handle motion을 현재 frame convention으로 변환한다.
  \item[\todo] Wind 값 변경만으로 basis, patch, binding을 재생성하지 않는다.
  \item[\todo] Frame runtime에서 material/attachment layout 변경 요청을 거부하고 새 setup/package version을 요구한다.
\end{itemize}

\subsection*{Module 2: current surface reconstruction}

\begin{itemize}
  \item[\todo] \(x=x^0+\Phi q+\sum_r\Psi_r z_r\), \(v=\Phi\dot q+\sum_r\Psi_r\dot z_r\)를 한 번 lift한다.
  \item[\todo] Raw rest/current local Jacobian \(\widehat J_s^0,\widehat J_s^t\in\mathbb R^{3\times2}\)와
  fixed rest-chart normalization
  \(H_s=[(\widehat J_s^0)^\top\widehat J_s^0]^{-1/2}\)를 만든다.
  \(J_s^0=\widehat J_s^0H_s\), \(J_s^t=\widehat J_s^tH_s=[a_{s,1},a_{s,2}]\)로 두어
  \((J_s^0)^\top J_s^0\approx I_2\)를 검사한다.
  \item[\todo] Rank-2 tangent affine map을
  \(F_{\parallel,s}=J_s^t(J_s^0)^\dagger\in\mathbb R^{3\times3}\)로 정의하되,
  area-normal은 이 3\(\times\)3 map의 cofactor가 아니라 \(J_s^t\)의 두 3D column을 직접 사용한다.
  \item[\todo] Area-normal을 3\(\times\)3 cofactor가 아니라
  \(\widetilde n_s=A_s^0(a_{s,1}\times a_{s,2})\)로 계산한다.
  비정규화 chart를 debug backend로 쓸 때만 대응 rest cross norm으로 나눈다.
  \item[\todo] \(A_s=\|\widetilde n_s\|\), \(n_s=\widetilde n_s/(A_s+\varepsilon)\)와
  \(\bar n_s^0=\operatorname{normalize}((J_s^0e_1)\times(J_s^0e_2))\)의 orientation continuity를 구현한다.
  \item[\todo] Gaussian covariance transport용 full shell map은 cross-product normal을 먼저 구한 뒤
  \(F_{\mathrm{shell},s}=F_{\parallel,s}+\lambda_{n,s}n_s(\bar n_s^0)^\top\in\mathbb R^{3\times3}\)로 별도 계산한다.
  \(\lambda_{n,s}\)의 default thickness/normal-extension 규칙을 package에 명시한다.
  \item[\todo] \code{SurfaceTangentMap3x2}=\(J_s^t\), \code{TangentAffineMapRank2\_3x3}=\(F_{\parallel,s}\),
  \code{FullShellMap3x3}=\(F_{\mathrm{shell},s}\)를 서로 다른 typed field로 두고 alias하지 않는다.
  \item[\todo] Tangent singular value/area clamp, orientation continuity와 degeneracy flag를 구현한다.
  \item[\todo] Rest/rigid rotation/in-plane scale/shear fixture로 normal과 area를 검증한다.
  \item[\todo] Aero normal과 rendering transport normal은 같은 orientation convention을 쓰되 서로 다른 map type을 올바르게 소비하는지 검사한다.
\end{itemize}

\subsection*{Module 3: relative wind와 analytic surface force}

\begin{itemize}
  \item[\todo] Spatial wind를 current sample position에서 query한다.
  \item[\todo] Relative wind를 \(W(x,t)-v\)로 정의한다.
  \item[\todo] Two-sided thin-sheet drag, friction, optional lift term을 구현한다.
  \item[\todo] Current normal/area와 material/aero coefficient를 사용한다.
  \item[\todo] Zero-relative-wind에서 force가 정확히 0인지 검사한다.
  \item[\todo] Uniform rigid translation에서 air/object velocity cancellation을 검사한다.
  \item[\todo] Front/back flip과 normal sign에 대한 force continuity를 검사한다.
  \item[\todo] Component별 force와 net force/torque를 debug output으로 남긴다.
\end{itemize}

\subsection*{권장 API}

\begin{lstlisting}[language=Python]
validate_and_prepare_state(state, package, handle_motion, dt)
    -> PreparedFrameState

reconstruct_surface(prepared_state, package)
    -> SurfaceState

evaluate_quasi_steady_aero(surface, wind_query, aero_parameters, time)
    -> AnchorAeroForces
\end{lstlisting}

\subsection*{필수 비교 backend}

\begin{itemize}
  \item Full current-surface/relative-wind reference.
  \item Rest normal/area ablation.
  \item Absolute wind instead of relative wind ablation.
  \item Object-average wind instead of spatial query ablation.
\end{itemize}

\textbf{완료 기준.}
Full-anchor evaluator가 analytic unit tests와 teacher surface snapshot의 force/net-wrench comparison을 통과하고,
dynamics를 적분하지 않는 isolated aero report가 생성된다.

\section{TD05. Always-on Global reduced physics predictor}

\textbf{목표.}
Oracle topology에서 Local을 끈 채 object-wide low-frequency wind response를 안정적으로 적분한다.

\textbf{선행 조건.} TD03 reduced operators, TD04 full current-surface aero.

\subsection*{구현 체크}

\begin{itemize}
  \item[\todo] Full-anchor force를 fixed \(\Phi\)로 projection하여 Global generalized load를 만든다.
  Optional H sample path도 setup-fixed \(T_s\)와 \(\Phi_s=T_s\Phi\)만 사용한다.
  \item[\todo] \(K_{gg}^0=\Phi^\top K^0\Phi\),
  \(C_{gg}=\Phi^\top C^0\Phi\)와
  \(Q_{g,\mathrm{op}}=\Phi^\top(h_K+h_C+f_{\mathrm p})\)를 분리한
  projective Global equation을 구현한다. Canonical source는
  \code{eq:global-operator-target-load}와 \code{eq:global-dynamics}다.
  \item[\todo] Canonical integrator는 \code{eq:global-implicit-midpoint}의
  implicit midpoint로 고정한다. Backward Euler는 robustness/numerical-damping,
  rest-linear \((K_{\mathrm{lin}}^0,C_{\mathrm{lin}}^0)\)는 structural closure ablation으로만 실행하며
  canonical result와 섞지 않는다.
  \item[\todo] Gravity, compliant attachment/handle, base aero, projective target/lift와
  structural cross term을 서로 다른 force-ledger channel로 기록한다.
  \item[\done] 이전 Local state의 cross term은 predictor에서 소비하고, TD11 corrector는 updated-minus-consumed delta만 소비하도록 계약을 고정했다.
  \item[\todo] Mainline의 optional learned Global generalized-force head는 default off로 둔다.
  \item[\todo] Mode rank 8/16/32/64 sweep와 factorization cache를 구현한다.
  \item[\todo] Setup material/attachment 변경이 새 package와 cache invalidation을 요구하는지 test하고, frame wind/handle motion은 basis를 invalidate하지 않는지 검사한다.
  \item[\todo] Zero wind/zero displacement rest invariance를 검사한다.
  \item[\todo] Wind start/stop에서 자연스러운 acceleration과 damping envelope를 검사한다.
  \item[\todo] Sinusoidal/chirp에서 modal frequency, amplitude, phase를 teacher와 비교한다.
  \item[\todo] Long rollout에서 energy growth/divergence를 검사한다.
\end{itemize}

\subsection*{권장 API}

\begin{lstlisting}[language=Python]
project_full_aero_to_global(aero_forces, surface, operators)
    -> ReducedAeroLoad

solve_global_predictor(
    previous_state, reduced_load, handle_motion, operators, dt
) -> GlobalPrediction
\end{lstlisting}

\subsection*{필수 baseline}

\begin{itemize}
  \item High-resolution teacher.
  \item Same-anchor sparse shell/PD/XPBD.
  \item WPB-style rest-state/uniform-wind modal projection.
  \item Direct MLP/GNN next-state Global.
  \item Recurrent Global state-space network.
  \item Physics Global + optional learned generalized-force residual.
\end{itemize}

Gate C의 K13 same-anchor sub-gate는 먼저 asset에 가장 잘 맞고
최적화된 same-anchor solver 하나로 판정한다.
최종 Gate C는 TD10 이후 K16/K18의 actual-skip overhead와 decay accumulation까지 함께 판정한다.
다른 PD/XPBD/corotational variant는 distinct Pareto 또는 failure mode가 있을 때 확장한다.
Teacher는 accuracy oracle이며 speed competitor가 아니다.

\subsection*{필수 지표}

Trajectory/velocity error, tip amplitude/phase, dominant frequency, damping ratio,
modal coordinate error, energy drift/divergence, p50/p95 dynamics-only latency와 peak memory.
Anchor 수 512/1024/2048, Gaussian 수와 multi-object batch를 독립 sweep하여
simulation resolution과 rendering resolution의 분리를 확인한다.

\textbf{완료 기준.}
MVP-A flag/strip에서 frequency와 damping이 teacher/same-anchor baseline과 일치하고,
4--10초 start/stop/chirp rollout이 발산하지 않는다.
Local 구현 전 mode-count Pareto를 저장하여 Global rank 증가만으로 충분한지 확인한다.
또한 same-anchor sparse shell/PD/XPBD 대비 개발 자원 배분용 내부 기준인
dynamics-only 약 2배 이상의 명확한 이득을 보이거나,
동일 p95 latency에서 velocity/trajectory/spectrum error를 유의미하게 줄이거나,
multi-object throughput 이득을 보여야 speed contribution을 유지한다.
세 경로가 모두 실패하면 K13에 따라 MPM 대비 속도만으로 main speed claim을 만들지 않는다.

\section{TD06. 조건부 H: current-state aerodynamic hyper-reduction}

\textbf{목표.}
TD04--TD05 profiler에서 full-anchor current-surface aero가 dynamics 비용의 유의미한 비율일 때만,
Global generalized load 정확도를 우선 유지하며 wall-clock를 줄일 수 있는지 검증한다.

\textbf{선행 조건.} TD04 full evaluator, TD05 Global basis/Jacobian과 rollout reference.

\begin{tcolorbox}[colback=LightOrange,colframe=orange!60!black,title=H 진입/철회 계약,fonttitle=\bfseries]
Full-anchor analytic evaluator가 V0 default이자 항상 사용 가능한 fallback이다.
Full-anchor가 병목이 아니면 TD06 구현을 중단해도 TD07 이후 core path는 계속 진행한다.
Attached object의 primary target은 \(Q_g=\Phi^\top f\) fidelity이며,
net force/torque는 auxiliary constraint/diagnostic이다.
\end{tcolorbox}

\subsection*{Offline snapshot와 cubature}

\begin{itemize}
  \item[\todo] Training deformation, velocity, spatial wind snapshot에서 full-anchor force를 계산한다.
  \item[\todo] 각 snapshot의 Global generalized force, net force, net torque를 label로 저장한다.
  \item[\todo] Force/torque center를 rest center, current center 중 하나로 고정하고 manifest에 기록한다.
  \item[\todo] Candidate sample support가 orientation, boundary, attachment, high-curvature 영역을 포함하는지 검사한다.
  \item[\todo] Nonnegative empirical cubature sample과 weight를 최적화한다.
  \item[\todo] Global generalized force를 primary objective로, net wrench를 auxiliary objective/constraint로 사용한다.
  \item[\todo] Held-out deformation과 traveling gust snapshot으로 overfit을 검사한다.
  \item[\todo] Sample 수별 실제 GPU kernel/gather/project latency를 측정한다.
  \item[\todo] H가 채택된 경우에만 estimated cubature error/confidence를 runtime diagnostic과 optional G feature로 전달한다.
\end{itemize}

\subsection*{Runtime Module 4}

\begin{itemize}
  \item[\todo] Current selected sample의 surface state와 relative wind만 평가한다.
  \item[\todo] Setup-fixed sample extraction \(T_s\)와
  \(\Phi_s:=T_s\Phi\)를 사용해 \(Q_g^{\mathrm{aero}}=\Phi_s^\top f_s\)를 만든다.
  Current state는 sample position/velocity/normal/area/relative wind와 force 값만 바꾸며
  Global lifting/Jacobian은 바꾸지 않는다. Canonical source는
  \code{eq:fixed-global-sample-jacobian}, \code{eq:full-generalized-aero}다.
  \item[\todo] State-dependent sample Jacobian이나 nonlinear sample map을 시험하면
  별도 versioned optional backend로 분리하고 build/update/corrector 비용을 profiler에 전부 포함한다.
  \item[\todo] Weighted net force와 net torque를 동시에 계산한다.
  \item[\todo] Full-anchor debug/reference backend를 같은 API로 제공한다.
  \item[\todo] Sample identity는 setup 후 고정하며 wind direction마다 anchor를 다시 고르지 않는다.
  \item[\todo] Weight의 음수/폭주, invalid sample, degeneracy를 runtime에서 검사한다.
\end{itemize}

\subsection*{권장 API}

\begin{lstlisting}[language=Python]
fit_aero_cubature(full_snapshots, operators, constraints, config)
    -> AeroCubature

reduce_aero_load(sample_forces, surface, aero_cubature)
    -> ReducedAeroLoad
\end{lstlisting}

\subsection*{E2 비교 backend}

\begin{enumerate}
  \item Full current-surface samples.
  \item Rest-state/uniform-wind WPB-style projection.
  \item Random samples.
  \item FPS samples.
  \item Learned samples without conservation.
  \item Canonical directional RSH/modal cache.
  \item Low-rank state-conditioned RSH.
  \item Direct direction-conditioned MLP/generalized-force backend.
  \item Proposed conservative current-state cubature.
\end{enumerate}

\begin{tcolorbox}[colback=LightOrange,colframe=orange!60!black,title=RSH 결정 규칙,fonttitle=\bfseries]
RSH는 mainline dependency가 아니다. 동일 sample/data/latency budget에서 held-out direction generalization이나
실제 wall-clock 이득을 명확히 제공할 때만 optional backend로 유지한다. 그렇지 않으면 harmonic direction encoding 또는 direct analytic evaluator를 사용한다.
\end{tcolorbox}

\textbf{완료 기준.}
Held-out deformation/spatial gust에서 full evaluator 대비 generalized-force와 net-wrench tolerance를 만족하고,
random/FPS/nonconservative backend보다 quality--latency 또는 drift Pareto가 좋다.
Full-anchor 대비 실제 wall-clock 개선이 없거나 generalized-load tolerance를 깨면 K14에 따라 H를 제거한다.
통과해도 H는 첫 논문의 supporting module이며 core contribution으로 자동 승격하지 않는다.

\section{TD07. Complementary Local physical system과 fixed-\(K\) Oracle active set}

\textbf{목표.}
Global rank가 놓치는 local flutter/fold를 network 없이 먼저 검증하고,
Global complement에서 하나의 Local physical state가 안정적으로 적분되는지 확인한다.

\textbf{선행 조건.}
TD07-A single-patch unit solver는 TD03 \(\Psi\)/Local operators와 TD05 Global predictor를 요구한다.
공식 TD07-B Oracle Local rollout은 TD11-A joint assembly/complement까지 통과한 force를 요구한다.

\subsection*{첫 구현 범위}

\begin{itemize}
  \item[\todo] 한 개 non-overlap patch와 patch-local fixed slot \(\Psi_r\)로 가장 작은 Local solver를 만든다.
  \item[\todo] TD11-A 통과 뒤 strongly overlapping two-patch fixture를 추가한다.
  \item[\todo] Oracle future error-reduction score로 fixed-\(K\) patch/overlap group을 선택한다.
  \item[\todo] Physics-only Local과 full-resolution Local upper bound를 만든다.
  \item[\todo] Network force는 0으로 두고 analytic/base local load만 사용한다.
\end{itemize}

\subsection*{Complement와 Local base}

\begin{itemize}
  \item[\todo] Force-space complement와 state-space mass orthogonality를 별도 API로 구분한다.
  \item[\todo] Local force proposal을 patch-indexed \(\Psi_r\) coordinate로 projection하기 전 Global leakage를 측정한다.
  \item[\todo] \(K_{\mathcal S_t\mathcal S_t}^0\),
  \(C_{\mathcal S_t\mathcal S_t}\)와
  \(Q_{\mathcal S_t,\mathrm{op}}=\Psi_{\mathcal S_t}^\top(h_K+h_C+f_{\mathrm p})\)를
  분리한 \code{eq:local-dynamics}의 projective Local equation을 구현한다.
  \item[\todo] Global state에서 오는 base aero/external/cross load를 정확히 한 번 넣는다.
  \item[\todo] Selector/activation 값이 \(M,C,K\)를 scale하지 않도록 한다.
  \item[\todo] Excitation blend \(\gamma_r\)는 TD11 Module 10의 force assembly에서 proposal에 정확히 한 번만 적용한다. Module 7, complement, Local solve, corrector에서 다시 곱하지 않는다.
  \item[\todo] \(S_t=A_t\cup D_t\)의 principal \(M_{rs},C_{rs},K_{rs}\) block을 모아 하나의 coupled Local coordinate system으로 적분한다.
  Active patch는 analytic excitation과 채택된 optional F를 제안하고, Decay patch는 F를 끄고 남은 state의 physical decay solve를 계속한다.
  \item[\todo] Persistent \(z_{\mathrm{sup}},\dot z_{\mathrm{sup}}\)에서 \(S_t\) fixed slots를 gather--solve--scatter하고 Dormant slot 처리와 decay 비용을 profile한다.
  \item[\todo] Canonical Local update는 \code{eq:local-effective-system}의 implicit midpoint exact
  active-principal-block solve로 고정한다. Backward Euler/rest-linear은 Step 0 ablation이고
  active-set factor cache key에는 package/operator version, integrator, \(\Delta t\)와 support signature를 넣는다.
\end{itemize}

\subsection*{V0 이후 dynamic basis 연구}

\begin{itemize}
  \item V0에서는 dynamic basis recomputation을 구현하지 않는다.
  \item 후속 최적화에서만 old/new basis state transfer, kinetic/potential energy jump와 coordinate initialization을 연구한다.
  \item Patch owner가 사라지거나 state-remap tolerance를 넘으면 fixed-superset으로 유지한다.
\end{itemize}

\subsection*{권장 API}

\begin{lstlisting}[language=Python]
solve_local_complement(
    previous_local_state, complement_force, global_prediction,
    operators, active_decay_block_indices, dt
) -> LocalPrediction
\end{lstlisting}

\subsection*{검증 체크}

\begin{itemize}
  \item[\todo] Patch별 \(\Phi^\top M\Psi_r\), inter-patch block rank/conditioning, force leakage와 modal amplitude inflation.
  \item[\todo] Zero-wind/selector-off energy monotonic decay.
  \item[\todo] Wind stop 후 physical recovery와 damping envelope.
  \item[\todo] Activation/deactivation position, velocity, acceleration, jerk continuity.
  \item[\todo] Same-p95 Global rank sweep 대 Oracle fixed-\(K\) Local의 velocity/trajectory/spectrum Pareto.
  Coordinate/rank-matched 비교는 보조 진단으로만 보고한다.
  \item[\todo] Patch-local normal, curvature, strain과 high-frequency PSD.
  \item[\todo] 4--10초 rollout과 timestep \(\times0.5,\times2\) sensitivity.
\end{itemize}

\textbf{완료 기준.}
Oracle Local이 동일 p95 dynamics latency의 Global rank sweep보다
velocity RMSE와 trajectory/PSD--spectrum을 반복적으로 개선하고,
zero-wind/activation/state transfer test에서 energy injection이나 discontinuity가 없다.
실패하면 K10에 따라 Local branch를 중단하거나 basis를 재설계한다.

\section{TD08. Teacher projection, 조건부 missing-force, benefit와 feedback label}

\textbf{목표.}
조건부 F가 위치 전체나 base force를 다시 학습하지 않도록 teacher trajectory에서 state와 missing force를 일관된 공간으로 변환하고,
fixed-\(K\) Oracle 분석과 조건부 G를 위한 future-benefit label 및 corrector label을 만든다.

\textbf{선행 조건.} TD02 teacher force channels, TD03 bases/operators, TD05/TD07 deterministic rollout.

\subsection*{Teacher state projection}

\begin{itemize}
  \item[\todo] Teacher surface state를 material anchors로 resample한다.
  \item[\todo] Mass projection으로 \(q^*,\dot q^*,\ddot q^*\)를 계산한다.
  \item[\todo] Global reconstruction을 제거한 complement에서 \(z^*,\dot z^*,\ddot z^*\)를 계산한다.
  \item[\todo] Teacher velocity/acceleration이 있으면 position finite difference보다 우선한다.
  \item[\todo] Position-only 데이터에는 smoothing, weak-form/integral consistency와 rollout supervision을 둔다.
  \item[\todo] Global/Local reconstruction residual과 patch complement energy를 저장한다.
\end{itemize}

\subsection*{Missing-force inverse dynamics}

\begin{itemize}
  \item[\todo] Local inertia, \(K^0,C^0\), frozen \(h_K,h_C\), hard lift \(f_{\mathrm p}\)와
  Global--Local cross block을 canonical runtime과 같은 timing/owner로 label 식에 포함한다.
  Rest-linear \(K_{\mathrm{lin}}^0,C_{\mathrm{lin}}^0\)를 canonical defect에 대신 넣지 않는다.
  \item[\todo] Gravity, base analytic aero, compliant target/operator contribution과 hard-support lift를
  force ledger에서 정확히 한 번 제거한다. Attachment를 generic external force로 다시 빼지 않는다.
  \item[\todo] Teacher solver가 직접 제공한 force breakdown을 우선 사용한다.
  \item[\todo] Structural, aerodynamic, attachment/reaction channel을 분리한다.
  \item[\todo] Teacher-FSI sample에만 unsteady aero defect를 표시한다.
  \item[\todo] Label noise/confidence와 inverse-dynamics balance residual을 저장한다.
  \item[\todo] Differentiable Local integrator의 multi-step rollout loss와 force label을 교차 검증한다.
\end{itemize}

\begin{tcolorbox}[colback=LightRed,colframe=BlockRed,title=조건부 F 전에 반드시 고정할 generalized-to-nodal 계약,fonttitle=\bfseries]
문서의 inverse dynamics는 Local generalized force \(\Delta Q_\ell^*\)를 정의하지만,
overlap/conservation pipeline은 patch nodal-force proposal \(\Delta f_r^*\)를 요구한다.
\textbf{기본 경로는 신뢰할 수 있는 channel별 teacher nodal-force breakdown}이다.
Canonical contract는 \code{eq:teacher-anchor-force-label}--style transfer가 아니라
각 stable label
\code{eq:teacher-anchor-force-label}, \code{eq:teacher-hard-free-force-split},
\code{eq:teacher-local-force-label},
\code{eq:teacher-hard-structural-augmented-label},
\code{eq:patch-force-label}을 따른다.
\begin{enumerate}
  \item Virtual-work conservative mesh-to-anchor transfer 뒤 known/base owner를 channel별로 한 번 뺀다.
  Hard layout에서는 free-support field와 raw teacher hard-row reaction을 먼저 분리한다.
  Hard-row 값은 patch Local head의 target이 아니다.
  \item External/free-object channel은 Global component를 먼저 route하고 남은 resultant와
  \(\Phi^\top f=0\)를 하나의 joint label solve에서 만족시킨다.
  Hard missing-structural channel은 runtime과 동일한 augmented variable
  \((y_{\mathrm{struct}}^*,\rho_{h,\mathrm{struct}}^*)\)와 constraint를 사용한다.
  \item F head의 supervised output은 patch별 \emph{raw free-support remainder}다.
  Owner-weighted reassembly 이후의 free force \(y^*\)와 hard companion \(\rho_h^*\)는
  post-assembly supervision/diagnostic이다. Head가 reaction을 직접 출력하거나
  raw/assembled target을 같은 tensor 이름으로 alias하지 않는다.
  \item Package-versioned \(W_f\), support-restricted \(W_{f,t}\), accepted-relation graph의 \(L_f\),
  smoothing coefficient와 one-sided attachment/boundary stencil을 label metadata에 저장한다.
  \(L_f\)가 hard row, attachment boundary 또는 rejected close-layer relation을 가로지르지 않게 한다.
\end{enumerate}

\textbf{Generalized-force lifting은 동등한 기본 경로가 아니다.}
Channel별 \(\Delta Q_{\ell,c}^{*,\mathrm{input}}\)와 독립적인 measured/teacher
conservation target \(b_{\mathrm{cons},c}^*\)가 모두 있는 free/no-hard sample에서만
\code{eq:generalized-force-lifting}의 optional 후속 경로를 허용한다.
\(b_{\mathrm{cons},c}^*\)가 없는 generalized-only sample은 사용하지 않는다.
Hard missing-structural generalized fallback은 raw support reference와 teacher reaction까지 있어야
동일 augmented objective를 복원할 수 있으므로 첫 논문에서는 no-go다.
Channel별 generalized split consistency가 실패하거나 aerodynamic/structural source가 불명확하면
single total-defect head로 합치지 않고 conditional F를 거부한다.
Transfer/lift/package version, constraint rank/feasibility와 raw--assembled reconstruction residual 없이
TD09를 학습하면 \killmark 처리한다.
\end{tcolorbox}

\subsection*{Counterfactual future-benefit label}

\begin{itemize}
  \item[\todo] 같은 initial state/wind에서 Global-only horizon rollout을 만든다.
  \item[\todo] 각 patch 또는 strongly overlapping group에 Oracle Local을 켠 counterfactual rollout을 만든다.
  \item[\todo] Position, velocity, normal, curvature, strain과 spectrum error 감소를 계산한다.
  \item[\todo] V0 label은 \(b_r=[\mathcal E_H^G-\mathcal E_H^{G+L_r}]_+\)인 positive error reduction으로 만들고 cost로 나누지 않는다.
  \item[\todo] Strongly overlapping patch는 group counterfactual label도 만들고 benefit이 patch별로 가산적이라고 가정하지 않는다.
  \item[\todo] Measured execution cost는 label 분모가 아니라 optional G/H profiling metadata로 별도 저장한다.
  \item[\todo] Short-horizon label과 일부 long-horizon calibration sample을 만든다.
  \item[\todo] Oracle active set, benefit rank, patch group과 horizon을 저장한다.
  \item[\todo] Teacher future state/force는 label에만 쓰고 runtime feature에서는 제거한다.
\end{itemize}

\subsection*{Feedback label}

\begin{itemize}
  \item[\todo] Predictor geometry의 Global aerodynamic load를 저장한다.
  \item[\todo] Oracle/corrected Local geometry의 Global aerodynamic load를 저장한다.
  \item[\todo] 두 load의 차이 \(\Delta Q_g^{\mathrm{aero},*}\)를 저장한다.
  \item[\todo] 실제 nonzero cross block의 updated-vs-consumed structural delta를 저장한다.
  \item[\todo] Same-frame, next-frame, no-feedback reference를 구분한다.
\end{itemize}

\subsection*{필수 산출물}

\begin{lstlisting}
projected_states/
  q_star, qdot_star, qddot_star
  z_star, zdot_star, zddot_star
  projection_error, complement_energy

optional_F_missing_force_labels/
  raw_free_support_patch_force_channels
  assembled_free_force_and_hard_companion
  local_generalized_force_and_channel_split_residual
  base_channels_removed
  W_f_W_f_t_L_f_and_boundary_versions
  constraint_rank_feasibility_confidence

optional_G_benefit_labels/
  positive_benefit, oracle_topk, optional_measured_cost
  error_components, horizon, patch_group

feedback_labels/
  predictor_load, corrected_load, aero_delta
  structural_consumed, structural_updated, structural_delta
\end{lstlisting}

\textbf{완료 기준.}
Label에서 base force를 다시 더하면 teacher balance가 깨지는 negative test와,
runtime-observable feature만으로 dataset을 읽는 leakage test가 통과한다.
Nodal-breakdown 기본 경로와 허용된 optional generalized lift가 구분되고,
raw-head 대 assembled \((y^*,\rho_h^*)\) target, hard-support no-go와
feedback consumed-force ledger가 문서와 코드에 고정된다.

\section{TD09. 조건부 F: always-on active-patch missing-force expert}

\textbf{목표.}
V0 physics-only Local을 고정한 뒤 모든 oracle-relevant patch에서 missing-force expert를 실행하여
network 자체가 추가로 필요한지 검증한다. 실패해도 TD10의 fixed-\(K\) physics selector와 core pipeline은 유지한다.

\textbf{선행 조건.} TD07 stable Local physics, TD08 force target 결정/labels.

\subsection*{입력 feature}

\begin{itemize}
  \item[\todo] Global prediction \(q,\dot q\)와 local history \(z,\dot z\).
  \item[\todo] Patch current position/frame, normal, area, strain, curvature와 deformation-rate invariant.
  \item[\todo] Spatial/relative wind와 base aero force.
  \item[\todo] Setup-fixed material/attachment metadata와 topology confidence.
  \item[\todo] Basis truncation indicator, 조건부 H를 쓸 때만 cubature error, optional recurrent memory.
  \item[\todo] 모든 vector/tensor를 canonical 또는 co-rotated patch frame으로 표현한다.
\end{itemize}

\subsection*{출력 계약}

\begin{itemize}
  \item[\todo] Stage 7은 Module 3--4가 이미 만든 owner-resolved
  \(f_{\mathrm{cache,ext}}\)를 active/decay support에 gather한다.
  Selector/F를 위해 wind query나 analytic aero를 다시 평가하지 않는다.
  \item[\todo] \code{PatchForceBatch}에서 analytic base aero, missing aero, missing structural을 별도 channel/owner/record ID로 유지한다.
  \item[\todo] Position, displacement 또는 next state가 아니라 missing force proposal을 출력한다.
  \item[\todo] F는 active patch의 raw free-support remainder만 출력한다.
  Hard-row reaction과 post-assembly \((y,\rho_h)\)를 head output에 넣지 않는다.
  \item[\todo] Missing aerodynamic external과 missing structural channel을 분리한다.
  Source가 불명확한 single total-defect head와 attachment/reaction head는 first-paper no-go다.
  \item[\todo] External resultant와 free/hard structural reaction semantics를 metadata에 둔다.
  \item[\todo] Confidence/uncertainty와 optional recurrent memory를 출력한다.
  \item[\todo] TD08의 nodal-breakdown 기본 경로 또는 모든 auxiliary conservation target을 갖춘
  optional free/no-hard lift만 사용한다. Runtime 출력은 항상 raw free-support patch nodal-force proposal로 고정한다.
  \item[\todo] Proposal 단계에서는 적분, overlap 평균, Global projection을 수행하지 않는다.
\end{itemize}

\subsection*{Model과 학습}

\begin{itemize}
  \item[\todo] 가장 작은 MLP/patch graph model을 baseline으로 만든다.
  \item[\todo] Temporal memory 없는 model을 먼저 학습한다.
  \item[\todo] Gust/history subset에서만 GRU/state-space memory의 추가 이득을 평가한다.
  \item[\todo] Force regression, direction/magnitude, channel consistency loss를 구현한다.
  \item[\todo] Differentiable Local solve를 통과한 one-step/rollout state loss를 추가한다.
  \item[\todo] Work/energy injection penalty와 confidence-weighted noisy-label loss를 평가한다.
  \item[\todo] Teacher forcing뿐 아니라 closed-loop predicted-state rollout을 curriculum에 넣는다.
  \item[\todo] Training normalization과 material/wind scale extrapolation을 기록한다.
\end{itemize}

\subsection*{권장 API}

\begin{lstlisting}[language=Python]
predict_patch_missing_forces(
    active_patch_batch, global_prediction, previous_local_state,
    wind_features, setup_metadata, model
) -> PatchForceBatch
\end{lstlisting}

\subsection*{필수 비교}

\begin{enumerate}
  \item Global-only.
  \item Full-resolution Local physics upper bound.
  \item Reduced Local physics-only.
  \item Reduced Local + learned missing force.
  \item Direct local displacement network.
  \item Learned force without Local physical integration.
  \item PGRD-style all-node/all-patch residual.
\end{enumerate}

\subsection*{검증}

Patch position/normal/curvature, flutter amplitude/frequency, high-frequency PSD,
wind-stop phase/decay, rollout stability, energy growth와 Local expert latency를 측정한다.

\textbf{완료 기준.}
Always-on learned missing force가 동일 Local physics-only보다 여러 seed/trajectory에서 반복적으로 개선하며,
direct displacement보다 wind-stop recovery와 long rollout이 안정적이다.
개선이 없으면 K11에 따라 F를 제거하고 physics-only Local을 유지한다.
F가 통과해도 learned G는 별도 TD10 비교를 통과해야 한다.

\section{TD10. V0 fixed-\(K\) selector, 조건부 G와 patch state machine}

\textbf{목표.}
먼저 학습이나 patch-cost 가산성 가정 없이 fixed cardinality에서 Local physics를 선택하고,
조건부 G가 runtime-observable feature로 같은 atomic grouping/\(K_A\)의 selection quality와
measured-latency Pareto를 실제로 개선하는지 별도 검증한다.
F를 채택한 경우 inactive F expert evaluation도 실제로 건너뛰되, F가 없어도 Local physics selector는 동작해야 한다.

\textbf{선행 조건.} TD11-A를 거친 TD07-B stable fixed-superset Oracle Local physics.
조건부 G에는 TD08 future-benefit labels가 추가로 필요하며, TD09는 F를 함께 dispatch할 때만 필요하다.

\subsection*{V0 selector와 고정 cardinality}

Module 3--4가 이미 계산한 free-DOF/owner-resolved anchor force를 재사용한다.
Every-frame selector를 위해 auxiliary rollout, 두 번째 geometry reconstruction 또는
aero/internal-force evaluation을 실행하지 않는다.
\(f_{\mathrm{cache,ext}}^t=P_{\mathrm{free}}
  (f_{\mathrm{aero}}^{\mathrm{base},t}+f_{\mathrm{gravity}}+f_{\mathrm{other,free}}^t
  )\),
\(f_{\mathrm{cache,sel}}^t=f_{\mathrm{cache,ext}}^t
  +P_{\mathrm{free}}(h_K^t+h_C^t+f_{\mathrm p}^t)\)로 둔다.
Compliant inhomogeneous target은 canonical projective target ledger의 \(h_K,h_C\)가 소유하며
별도 attachment RHS로 다시 더하지 않는다.
Atomic gating group \(\mathcal G\)의 cached driving load, fixed effective operator와
owner-mass-normalized score는 canonical
\code{eq:cached-local-drive}, \code{eq:cheap-dynamic-compliance}를 그대로 소비한다.
Schema assertion은 \(K^0,C^0\), implicit-midpoint 계수
\(a_M=4\Delta t^{-2},a_C=2\Delta t^{-1}\), mass-valued \(\varepsilon_m\)과
setup/fixed-timestep prefactorization을 검사한다. Backward-Euler 계수는 diagnostic ablation type에만
허용한다. 이 checklist에 표시한 기호와 계수는 독립 normative equation이 아니며,
TD01 canonical-reference test가 위 label과 schema assertion의 drift를 막는다.
\(m_{\mathcal G}^{\mathrm{own}}\)은 partition-of-unity owner weight로 각 anchor mass를
group 안에서 한 번만 합산한 값이다.
Prescribed hard support는 \(f_{\mathrm p}\), compliant target은 \(h_K,h_C\)에만 들어가며
selector가 별도 attachment force로 다시 만들지 않는다.
V0 active set은 Oracle benefit 또는 \(\eta_{\mathcal G,\mathrm{comp}}\)의
\(\operatorname{TopK}(K_A)\)이며 \(K_A\) sweep으로 quality--latency curve를 만든다.
Patch score의 합이나 patch별 cost가 overlap group의 joint benefit/wall-clock와 가산적이라고 가정하지 않는다.
Full Global-only counter-state residual은 offline 또는 \(J\)-frame low-rate analysis ceiling으로만 구현하고,
추가 solve/geometry/aero/internal evaluation과 activation delay를 모두 비용에 포함한다.

\begin{itemize}
  \item[\todo] Never, Always, matched-\(K\) Random, curvature, strain, every-frame cached-compliance와 Oracle future-benefit selector를 동일 \(K_A\)로 구현한다.
  \item[\todo] Full counter-state residual은 offline/low-rate analysis backend로만 구현하고 추가 solve/force evaluation과 delay를 포함한다.
  \item[\todo] \(K_A\in\{0,\ldots,K_{\max}\}\) sweep 또는 합리적인 subset으로 quality--latency curve를 만든다.
  \item[\todo] Strongly overlapping patch는 TD08에서 정의한 group 단위로 선택한다.
  \item[\todo] Patch별 benefit 합과 patch별 profiler cost 합을 joint set의 benefit/wall-clock로 간주하지 않는다.
  \item[\todo] Cross-group excitation, high-curvature/low-current-load release,
  topology/operator-error negative control과 fast traveling/reversing gust를 mandatory blind-spot fixture로 고정한다.
  \item[\todo] Pure compliance, dimensionless deterministic hybrid, low-rate/full counter와
  Oracle을 같은 fixed cardinality \(K_A\)와 measured p95에서 비교한다.
  Hybrid가 비용 포함 Pareto를 개선할 때만 runtime 후보로 승격하고, 실패하면 K24를 적용한다.
  \item[\todo] Teacher-residual threshold는 runtime 불가능한 analysis-only upper bound로 명확히 분리한다.
  \item[\todo] On/off threshold, hysteresis와 minimum dwell time을 구현한다.
  \item[\todo] Smooth excitation blend \(\gamma\)를 계산하되 Module 10 assembly에만 전달하고 Local \(M,C,K\), state 또는 다른 force stage에는 곱하지 않는다.
  \item[\todo] Inactive patch/group을 active principal-block gather와 조건부 F batch에서 제외한다.
  \item[\todo] Small active set에서 gather/scatter, factorization/solve와 kernel launch overhead의 break-even을 측정한다.
  \item[\todo] Always-on이 더 빠른 active ratio에서는 runtime dispatcher가 fallback할 수 있게 한다.
\end{itemize}

\subsection*{조건부 G: learned future-benefit selector}

\begin{itemize}
  \item[\todo] Pooled Global state와 cheap patch invariant만 사용하는 lightweight benefit ranker를 만든다.
  \item[\todo] Teacher future state, teacher residual/force를 runtime feature에서 제거한다.
  \item[\todo] G의 첫 버전 출력은 positive future-error reduction score와 uncertainty이며 predicted cost로 나누지 않는다.
  \item[\todo] Calibration loss/temperature calibration과 uncertainty threshold를 구현한다.
  \item[\todo] Cached-compliance와 low-rate/full-counter signal을 feature로 넣은/뺀 learned ranker ablation을 제공하고 추가 비용을 분리한다.
  \item[\todo] Gate feature construction, model inference와 selection 비용을 별도 profile한다.
  \item[\todo] Hardware-aware predicted-cost/TopBudget는 fixed-\(K\) 결과와 profiler가 필요성을 보인 뒤의 optional extension으로만 구현한다.
\end{itemize}

\subsection*{Active/Decay/Dormant state}

\begin{longtable}{L{0.16\linewidth}L{0.34\linewidth}L{0.42\linewidth}}
\toprule
상태 & 실행 & 전환 계약 \\
\midrule
Dormant & Cheap selector만 실행 & Fixed-\(K\) selection과 on 조건을 통과하면 Active. Local state는 충분히 작아야 함. \\
Active & Module 10에서 \(\gamma\)가 한 번 적용된 new excitation + Local physical solve; optional F & Off 조건과 dwell을 만족하면 Decay. State reset 금지. \\
Decay & Optional F off, Module 10 excitation fade, \(S=A\cup D\) Local physical decay solve & Owner/group energy와 selector score가 충분히 작을 때만 Dormant. 재상승하면 Active. \\
\bottomrule
\end{longtable}

Actual-skip trace는 Dormant의 Local force/solve/expert 전체 skip,
Decay의 새 expert/새 excitation skip과 mandatory physical decay solve를 서로 분리한다.
Decay solve를 끄거나 state를 reset해 speedup을 만드는 구현은 허용하지 않는다.
Patch/group fixed coordinate rank를 \(r_r\), core--halo unique-anchor support를
\(\mathcal I_r\), full superset rank를 \(r_{\mathrm{sup}}\)라 두고
\(r_A=\sum_{r\in\mathcal A_t}r_r\),
\(r_D=\sum_{r\in\mathcal D_t}r_r\),
\(r_S=r_A+r_D\)와 각각의 normalized ratio를 기록한다.
Unique support size는
\(n_A=|\bigcup_{r\in\mathcal A_t}\mathcal I_r|\),
\(n_D=|\bigcup_{r\in\mathcal D_t}\mathcal I_r|\),
\(n_S=|\bigcup_{r\in\mathcal A_t\cup\mathcal D_t}\mathcal I_r|\)로 별도 보고한다.
Patch/group cardinality \(|\mathcal A_t|,|\mathcal D_t|\), coordinate rank와
unique-anchor support size를 서로 대체하지 않는다.

\begin{itemize}
  \item[\todo] Active에서 Decay로 갈 때 \(z,\dot z\)를 보존한다.
  \item[\todo] Decay에서 physical energy가 비증가하는지 검사한다.
  \item[\todo] Recurrent memory의 decay/cleanup 정책을 정한다.
  \item[\todo] Toggle count, activation duration, acceleration/jerk를 기록한다.
  \item[\todo] Active/decay/dormant index의 장시간 memory leak를 검사한다.
  \item[\todo] Stationary/slow-traveling/fast-traveling/reversing gust별로
  \(r_A,r_D,r_S,n_A,n_D,n_S\)와 normalized rank ratio,
  decay residence distribution과 median/p95,
  unique-signature count/churn, cache hit/miss,
  gather/factor/solve/scatter p50/p95/p99와 always-on 대비 latency를 기록한다.
\end{itemize}

\subsection*{권장 API}

\begin{lstlisting}[language=Python]
select_local_patches(
    global_prediction, patch_features, previous_state,
    fixed_k, selector, optional_gate_model=None
) -> SelectorDecision

update_patch_state_machine(
    selector_decision, local_energy, previous_state, thresholds
) -> PatchStateUpdate
\end{lstlisting}

\subsection*{필수 비교와 지표}

Never, Always, matched-\(K\) Random, curvature, strain, cached-compliance,
low-rate/full-counter와 teacher-residual analysis-only, Oracle benefit 및 조건부 G learned benefit를 비교한다.
조건부 cost-aware extension은 별도 profiler ablation으로 둔다.
High-benefit recall/precision/AUPRC, calibration, Oracle regret, activation delay/high-benefit miss,
active/solve-set ratio,
실제 active Local solve 수와, F를 채택한 경우 expert call 수, kernel 수, p50/p95 latency, toggle/jerk를 보고한다.

\textbf{완료 기준.}
Core 완료는 fixed-\(K\) Oracle/cached-compliance selector, state decay와 실제 active-block skip이 재현되는 것이다.
G는 동일 \(K_A\)와 동일 p95 latency 두 조건에서 cached-compliance/heuristic보다 낮은 dynamics error를 보이고,
Always-on 대비 실제 wall-clock를 줄일 때만 채택한다. Classification score만 높은 것은 G 완료가 아니다.

\section{TD11. Conservative overlap, complement와 Local-to-Global corrector}

\textbf{목표.}
여러 active patch의 force를 하나의 물리계로 조립하고,
Global double counting을 제거한 Local update가 corrected geometry를 통해 Global state에 실제로 feedback되게 한다.

\textbf{선행 조건.}
TD11-A assembly/complement는 TD03 basis/operator와 TD04--TD05 force contract 뒤,
공식 TD07 Oracle Local Pareto 전에 synthetic/Oracle force fixture로 고정한다.
TD11-B corrected-aero/structural-delta corrector는 안정적인 TD07 rollout 뒤 실행한다.
TD10-core selector/state metadata는 adaptive trace에만 필요하다.
조건부 F 또는 G를 채택할 때만 각각 TD09 또는 TD10-G learned artifact를 추가로 요구한다.

\subsection*{실제 frame 내부 순서}

\begin{enumerate}
  \item Module 7이 active patch별 analytic Local base aero를 출력하고, 조건부 F를 채택한 경우에만 missing aero/structural proposal을 별도 channel로 출력한다.
  \item Module 10이 \(S_t=A_t\cup D_t\)의 각 channel을 anchor force space에 scatter하고 \(\gamma_r\)를 proposal에 정확히 한 번 적용한다. Decay patch의 optional F는 실행하지 않는다.
  \item Module 10--9가 overlap/conservation과 \(\Phi^\top f=0\)를 하나의 joint KKT에서 동시에 만족시킨다. 순차 projection은 허용하지 않는다.
  \item Module 8이 channel별 Local generalized force를 처음 합쳐 fixed-superset의 \(S_t\) principal block 하나에서 Local state \(z,\dot z\)를 한 번 적분한다.
  \item Module 11이 corrected aerodynamic delta와 updated-minus-consumed structural cross delta로 Global corrector를 한 번 푼다.
\end{enumerate}

\subsection*{Module 10: overlap-aware force assembly}

\begin{itemize}
  \item[\todo] Patch-to-global scatter와 core--halo partition weight를 구현한다.
  \item[\todo] Selector의 \(\gamma_r\)를 anchor-force proposal에 정확히 한 번 적용하고 force ledger에 \code{gamma\_consumed\_at=module10}을 남긴다.
  \item[\todo] Analytic base aero와 조건부 F의 missing aero/missing structural을 assembly 이후까지 서로 다른 owner/lineage로 유지한다.
  \item[\todo] Active/decay core--halo가 만지는 unique anchor support \(\mathcal I_t\)를 만들고,
  restricted force injection, Global row와 current-wrench operator를
  \code{eq:support-restricted-operators}에 맞춰 조립한다.
  \item[\todo] External aerodynamic residual은 같은 support 안에서 Global component를 먼저 route한 뒤
  Local remainder의 unique resultant를 보존한다. Full-domain Global field를 만든 뒤 잘라내거나
  support 밖 remainder를 Local KKT로 상쇄하지 않는다.
  \item[\todo] Free object의 intrinsic structural correction만 zero net wrench를 만족시킨다.
  Hard support에서는 free-node correction과 object-to-support companion reaction을 함께 닫고,
  compliant attachment는 structural/internal residual이 아닌 별도 owner로 유지한다.
  \item[\todo] Patch-copy force를 unique shared-anchor support로 먼저 condense한다.
  Owner partition은 scatter invariant로 검사하고, shared-anchor position/velocity continuity는
  persistent state와 single coupled solve로 보장한다. 이를 force-KKT row로 중복 추가하지 않는다.
  \item[\todo] Setup-defined block-diagonal SPD \(W_{f,t}\), support injection \(E_t\),
  channel별 \(B_c^t,d_c^t\)와 current-geometry force/torque operator를 구현한다.
  \item[\todo] Dense anchor-primal reference solver와 production unique-anchor small-Schur solver를 분리한다.
  \item[\todo] Independent constraint-row inventory, rank/condition과 channel별 per-RHS consistency를 검사한다.
  Infeasible/rank-deficient frame은 routing/constraint를 수정하거나 fail-fast한다.
  Aggregate feasible rate 또는 ridge/row deletion으로 특정 RHS 실패를 숨기지 않는다.
  \item[\todo] Fixed operator block과 geometry-dependent torque Gram update를 분리하고,
  exact signature가 같은 channel/frame에서만 factor를 공유한다.
  \item[\todo] Force/objective/constraint residual, Schur condition/pivot,
  raw-to-routed force adjustment, reaction amplification, generalized-work 보존,
  factor share/rebuild와 scatter/operator-build/factor/dual-solve/backprojection timing을 기록한다.
  Canonical diagnostics source는 \code{eq:support-routing-diagnostics},
  \code{eq:joint-feasibility-residual}, \code{eq:structural-schur-adjustment-diagnostics}다.
  \item[\todo] Development split에서 attachment class, support-size/rank,
  constraint-row inventory와 conditioning으로 나눈 signature class/bin별
  route/adjustment/reaction amplification, support expansion 횟수, 최대 halo 폭/anchor ratio와
  factor/cache budget을 fit/freeze한다. Test의 새 exact active-set ID는 검증된 class에 속하고
  runtime diagnostic을 통과하면 허용하되, unseen class나 frozen cap을 넘으면
  더 넓은 full-domain solve로 조용히 후퇴하지 않고 package를 reject하거나
  사전 선언한 always-on Local fallback과 최초 초과 시점을 기록한다.
  \item[\todo] Patch partition/core--halo 폭 변경에 대한 accuracy, feasible rate,
  halo-expansion frequency와 p95/p99 sensitivity를 검사한다.
  \item[\todo] 단순 displacement/force average는 ablation backend로만 제공한다.
\end{itemize}

\subsection*{Module 9: Global force complement}

\begin{itemize}
  \item[\todo] Force-space projector와 displacement/state projector의 mass metric을 명시한다.
  \item[\todo] Owner/overlap은 unique-anchor condensation으로 처리하고,
  internal action--reaction은 가능한 antisymmetric edge-force parameterization으로 내재화한다.
  남은 low-row resultant constraint와 \(\Phi^\top f_\ell^\perp=0\)만 하나의 joint Schur system에서 푼다.
  \item[\todo] Joint solution의 generalized leakage와 각 conservation residual tolerance를 검사한다.
  \item[\todo] First-paper method는 별도 Global defect head를 두지 않고 routed component를 diagnostic/loss로만 기록한다.
  \item[\todo] Corrected-aero delta와 legitimate cross reaction은 complement 제거 대상이 아님을 channel로 보장한다.
  \item[\todo] Wrench/conservation projection 뒤 complement를 순차 적용하거나 그 반대 순서로 적용하는 code path를 negative test로 금지한다.
  \item[\todo] Runtime fallback도 conservation과 \(\Phi^\top f=0\)를 동시에 만족해야 하며,
  torque 삭제나 sequential projection으로 후퇴하지 않는다.
\end{itemize}

\subsection*{Module 8: single Local update}

\begin{itemize}
  \item[\todo] Assembly/complement 결과를 channel별 \(Q_{\ell,c}^{\perp}\)로 변환한다.
  \item[\todo] Active와 Decay patch의 fixed slots를 모아 \(M_{rs},C_{rs},K_{rs}\) active principal block 하나에서 적분한다.
  \item[\todo] Base analytic Local aero와 사용 가능한 조건부 F defect channel은 여기서 처음 합치고, cross load는 별도 structural coupling 항으로 소비한다.
  \item[\todo] Patch별 solver/integrator 호출 후 결과를 average하는 code path가 없음을 trace test로 보장한다.
  \item[\todo] Patch self energy와 overlap-group/cross-block energy를 구분해 Local work, residual, state transition diagnostics와 함께 저장한다.
\end{itemize}

\subsection*{Free/hard/compliant reaction과 post-step residual}

\begin{itemize}
  \item[\todo] Free object는 intrinsic structural correction의 zero-wrench residual을 기록하되
  support reaction을 생성하지 않는다.
  \item[\todo] Hard support의 channel companion은
  \code{eq:hard-reaction-wrench-balance}의 object-to-support 부호를 사용한다.
  이는 support 밖의 정당한 reaction channel이며 Local excitation/complement force로 재투영하지 않는다.
  \item[\todo] Runtime total hard support-on-object reaction은 corrector와 final \(q,z\)가 확정된 뒤,
  \code{eq:post-step-midpoint-kinematics}의 discrete acceleration과
  \code{eq:full-applied-external-force-ledger}를 사용해 post-step diagnostic으로만 복원한다.
  \item[\todo] \code{eq:hard-total-reaction-residual}과
  \code{eq:hard-reaction-rom-residual-decomposition}에 따라
  total hard reaction, free-DOF \(r_{\mathrm{ROM}}\), Global/Local virtual-work residual을 분리한다.
  Hard reaction 하나로 reduced truncation residual까지 닫았다고 보고하지 않는다.
  \item[\todo] Compliant support는 \code{eq:compliant-reaction-ledger}의
  object force와 handle reaction을 분리하고 common-motion virtual-work balance를 검사한다.
  \item[\todo] Prescribed hard lift, compliant force, projective targets, base/corrected aero와
  optional F field가 post-step full-force ledger에 정확히 한 번만 들어가는 negative owner test를 만든다.
\end{itemize}

\subsection*{Module 11: feedback와 corrector}

\begin{itemize}
  \item[\todo] Predictor와 Local-corrected anchor/surface state를 각각 복원한다.
  \item[\todo] Active/halo region 또는 필요 sample에서 corrected aero를 재평가한다.
  \item[\todo] \(\Delta Q_g^{\mathrm{aero}}=Q_g^{\mathrm{corrected}}-Q_g^{\mathrm{predictor}}\)만 corrector에 넣는다.
  \item[\todo] Full base aero가 corrector에 다시 들어가면 실패하는 unit test를 만든다.
  \item[\todo] Nonzero \(K_{g\ell},C_{g\ell}\)일 때만 structural feedback을 사용한다.
  \item[\todo] One-corrector를 default로 구현한다.
  \item[\todo] Under-relaxation, force-delta clamp와 smaller timestep option을 둔다.
  \item[\todo] No/next-frame/same-frame one-corrector/multiple-iteration을 비교한다.
\end{itemize}

\begin{tcolorbox}[colback=LightRed,colframe=BlockRed,title=Structural feedback 소비 계약,fonttitle=\bfseries]
Module 5는 이전 Local state의 cross coupling을 소비하고 그 값을 \code{predictor\_cross\_load}로 기록한다.
Module 11은 updated reaction 전체가 아니라
\code{actual-midpoint coupling - predictor-midpoint coupling}인
\code{structural\_cross\_delta}만 더한다.
Canonical definition은 \code{eq:structural-feedback}이며,
actual midpoint가 predictor midpoint와 같으면 delta가 정확히 0이어야 한다.
Corrector는 endpoint 차이에 generic \code{SolveGlobal}을 호출하지 않고
\code{eq:global-incremental-corrector}를 사용한다.
TD01 schema assertion은 factor $2$, velocity update $2/\Delta t$와
implicit-midpoint operator coefficients가 canonical label과 일치하는지만 검사하며,
이 checklist에 별도 corrector equation을 복제하지 않는다.
Endpoint corrected-aero delta를 held midpoint-load delta로 적용하는 partitioned approximation이므로
monolithic/iterated midpoint reference와의 차이를 별도 보고한다.
Force ledger와 negative double-count test 없이 TD11을 완료 처리하지 않는다.
\end{tcolorbox}

\subsection*{권장 API}

\begin{lstlisting}[language=Python]
build_joint_constraint_operator(
    patch_force_batch, selector_decision, package, operators
) -> ConstraintOperator

solve_joint_constraints_schur(constraint_operator)
    -> AssemblyResult, ConstraintSolveReport

solve_joint_constraints_dense_reference(constraint_operator)
    -> AssemblyResult

validate_joint_constraint_equivalence(
    assembly_result, dense_reference
) -> ContractReport

solve_local_complement(...)
    -> LocalPrediction

correct_global_state(
    global_prediction, local_prediction, predictor_surface,
    wind_query, operators, consumed_force_ledger, dt
) -> CorrectedDynamicState

recover_support_reaction_ledger(
    corrected_state, previous_state, applied_force_ledger,
    package, dt
) -> SupportReactionLedger
\end{lstlisting}

\subsection*{필수 검증}

\begin{itemize}
  \item[\todo] Overlap boundary position/velocity/force jump.
  \item[\todo] Free object의 intrinsic structural net wrench와 object drift/rotation;
  hard support에서는 free-node force 단독이 아니라 companion reaction을 합친 sign-resolved wrench.
  \item[\todo] Global force/state leakage와 modal amplitude inflation.
  \item[\todo] Dense primal reference 대 small-Schur force/objective/constraint residual equivalence.
  \item[\todo] External route와 structural/hard-reaction augmented route 각각의 per-RHS feasible rate,
  adjustment/reaction amplification, generalized-work, halo expansion/cap과 final support ratio.
  \item[\todo] Constraint row 수가 overlap patch 수에 비례하지 않고,
  support/geometry/row signature 변경 시 factor가 정확히 invalidate되는지 검사한다.
  \item[\todo] Unique scatter, operator update, Schur build/factor/dual solve/backprojection의 p50/p95/p99와 dynamics latency share.
  \item[\todo] Corrector base-force/cross-force double-count negative test.
  \item[\todo] Free/hard/compliant fixture의 reaction 부호/wrench와 post-step
  \(r_{\mathrm{ROM}}\) 분해; synchronized reference와 one-corrector residual을 구분한다.
  \item[\todo] Local correction 전후 Global modal state와 teacher error 차이.
  \item[\todo] Corrected-aero delta magnitude/phase와 extra latency.
  \item[\todo] Strong follower-force, abrupt gust와 timestep stress test.
\end{itemize}

\textbf{완료 기준.}
TD11-A는 proposed assembly/complement가 dense reference와 일치하고 simple averaging보다 seam,
net wrench와 drift를 줄이는 것이다. TD11-B는 one-corrector가 base force를 중복하지 않으면서
적어도 일부 benchmark의 Global response/teacher error를 개선하는 것이다.
Feedback 효과가 없으면 closed-loop feedback SYS claim을 제거하고 Local-only 또는 next-frame update로 단순화한다.

\section{TD12. Static-GS topology/operator distillation과 predicted package}

\textbf{목표.}
Training-only source topology privilege를 이용하되, target setup에는 static 3DGS와 명시적 metric/material/attachment metadata만 받아
explicit target triangle-mesh reconstruction 없이 sparse oriented material scaffold와 reduced object package를 생성한다.

\textbf{선행 조건.}
Step 0 중에는 coordinate/schema, relation rank/valence, area/mass와 close-layer false edge를
보는 최소 readout-only probe만 허용한다.
Formal TD12-R은 Step-0 gate 통과 뒤 Track A0와 병렬로 시작한다.
Full TD12 학습은 TD03 Oracle package, TD12-R 통과와 TD05/TD07/TD10-core/TD11에서 downstream-sensitive operator/basis 요구를 확인한 뒤 시작한다.

\subsection*{TD12-R: 조기 representation feasibility kill track}

Adaptive dynamics 전체 구현 전에 5--10개의 simple thin-surface object와 각 object에서 독립적으로 생성한 여러 static 3DGS를 사용한다.

\begin{itemize}
  \item[\todo] Oracle surface와 independent GS마다 anchor/oriented point--edge--local-relation graph를 복원한다.
  \item[\todo] Total/local area와 user density/thickness로 유도한 mass 보존 오차를 측정한다.
  \item[\todo] 동일 setup material/attachment에서 첫 8--16개 eigenfrequency 오차를 비교한다.
  \item[\todo] 대응 modal subspace의 principal angle과 downstream Global projection 오차를 측정한다.
  \item[\todo] Modal/eigenfrequency error 대 downstream rollout error를 asset/load별 scatter로 저장하고 단순 단조관계로 과장하지 않는다.
  \item[\todo] 접힌 면, 얇은 간격의 이중 layer와 front/back에서 false edge/hinge shortcut을 계수한다.
  \item[\todo] Reconstruction seed, Gaussian density, hole/floater 변화에서 package validation과 실패 reason을 기록한다.
\end{itemize}

TD12-R이 반복적으로 area/mass, low modes 또는 close-layer topology를 보존하지 못하면,
full adaptive system을 기다리지 않고 topology distillation 범위를 줄이거나
``without explicit target triangle-mesh reconstruction'' 주장을 제거한다.

\subsection*{Surface reliability와 anchor selection}

\begin{itemize}
  \item[\todo] Opacity, covariance surface-likeness, multi-view contribution과 isolation feature를 계산한다.
  \item[\todo] Boundary/curvature/layer evidence와 reconstruction confidence를 추가한다.
  \item[\todo] Raw Gaussian count를 mass/area/importance로 직접 사용하지 않는다.
  \item[\todo] \(K\ll N_G\) material anchor를 coverage와 reliability 기준으로 선택한다.
  \item[\todo] Anchor rest normal/area weight와 geometry/reconstruction reliability를 예측한다.
  \item[\todo] User density/thickness와 area weight에서 mass를 유도하고 partition-of-unity로 total area/mass를 보존한다.
\end{itemize}

\subsection*{Oriented topology와 thin-shell relation}

\begin{itemize}
  \item[\todo] Geometry/evidence 기반 candidate edge generator를 만든다.
  \item[\todo] Material adjacency, rest length와 tangent를 예측하되 stretch/shear/bending 계수는 user material metadata와 geometry에서 조립한다.
  \item[\todo] Bending/hinge relation과 orientation consistency를 예측한다.
  \item[\todo] 가까운 disconnected layer와 front/back shortcut을 거부한다.
  \item[\todo] User attachment layout을 anchor constraint로 transfer하고 hard/compliant parameter를 그대로 보존한다.
  \item[\todo] Core--halo patches와 overlap graph를 생성한다.
  \item[\todo] Area-share logit과 별도로 membrane/bending relation-validity probability를 예측하고,
  reconstruction evidence와 MLS condition을 입력으로 object-disjoint development split에서 calibrate/freeze한다.
  Canonical source는 \code{eq:relation-validity-confidence}다.
  \item[\todo] Candidate count/valence에 의존하는 raw score sum을 absolute confidence로 쓰지 않고,
  count-aware anchor aggregator와 operator-constructibility target을
  \code{eq:relation-absolute-confidence}, \code{eq:anchor-absolute-confidence-target}에 맞춰 구현한다.
  \item[\todo] Relation validity, anchor absolute confidence, attachment-transfer residual과
  support-expand/fallback/reject reason을 channel별로 저장한다.
\end{itemize}

\subsection*{Operator/basis/package construction}

\begin{itemize}
  \item[\todo] Arbitrary dense matrix를 network가 직접 출력하지 않는다.
  Predicted \(\mathcal Q_{\mathrm{op}}\), calibrated validity/area share와 setup material에서
  \(M,K_{\mathrm{lin}}^0,C_{\mathrm{lin}}^0\), canonical
  \(K^0,C^0\)와 target descriptor ledger를 TD03과 같은 typed construction으로 조립한다.
  \item[\todo] Predicted package에서 TD01의 PSD/symmetry/null-space뿐 아니라
  rigid/equivariance, normal-curvature pollution, action/modal closure와 target-ledger test를 다시 실행한다.
  \item[\todo] Global basis와 patch-local fixed-superset Local basis/cross block을 재구성하고 TD01의 per-patch orthogonality/active-block conditioning을 검사한다.
  \item[\todo] Full-anchor aerodynamic evaluator를 predicted anchors에 구성하고, 조건부 H를 채택한 경우에만 sample/cubature를 transfer 또는 refit한다.
  \item[\todo] Static Gaussian-to-anchor affine binding을 생성한다.
  \item[\todo] Setup package가 explicit target triangle-mesh reconstruction, FEM/MPM particle grid를 참조하지 않는지 trace한다.
  \item[\todo] Package checksum, version, confidence와 runtime compatibility report를 저장한다.
\end{itemize}

\subsection*{권장 API}

\begin{lstlisting}[language=Python]
infer_anchor_scaffold(
    canonical_gs, setup_metadata: SetupMetadata, model, config
) -> AnchorScaffold

build_object_package(
    canonical_gs, anchor_scaffold, setup_metadata, setup_models, config
) -> ObjectPackage
\end{lstlisting}

\subsection*{Fallback와 정직한 scope}

\begin{itemize}
  \item[\todo] Low absolute confidence에서는 support 확대와 MLS rank/conditioning 재검사를 먼저 수행한다.
  여전히 실패하면 user hint, explicit conservative fallback 또는 package rejection을 반환하며,
  confidence를 stiffness/material scale에 곱해 실패를 숨기지 않는다.
  \item[\todo] 추가 geometry hint/category prior를 사용한 결과를 geometry-scaffold automatic 결과와 분리한다. 필수 user material/attachment metadata는 양쪽에 동일하게 제공한다.
  \item[\todo] Explicit target triangle-mesh reconstruction fallback은 별도 baseline/flag로 기록한다.
  \item[\todo] Stable area-normal/bending에 watertight triangle mesh가 사실상 필수라면 ``without explicit target triangle-mesh reconstruction'' claim을 제거하거나 범위를 명시적으로 약화한다.
\end{itemize}

\subsection*{E1 비교와 지표}

Oracle topology, proposed predicted topology, Euclidean kNN, FPS+kNN,
surface reconstruction mesh, no orientation, no bending/stretch separation,
user attachment를 oracle anchor에 직접 지정한 경우와 predicted anchor로 transfer한 경우를 비교한다.
Edge/hinge F1, layer shortcut, normal consistency, attachment-transfer error, area/mass error,
relation-validity Brier/ECE/AUPRC, anchor absolute-confidence coverage/false-accept rate,
modal frequency/shape와 downstream rollout/aero/stability를 함께 보고한다.
Calibration 지표는 object-disjoint test에서 anchor density, candidate count, valence와
condition-number bin별로 함께 보고한다. Topology/validity confidence는 triage 근거일 뿐이며,
area/mass, rank/condition, affine/quadratic reproduction, 초기 8--16 modes/modal angle와
common-probe rollout/stability hard check를 대체하지 않는다.

\textbf{완료 기준.}
Independent static-GS variants에서 package validation이 통과하고,
TD12-R의 area/mass, 첫 8--16 eigenfrequency, modal principal angle과 close-layer test가 통과하며,
Predicted topology의 downstream dynamics error가 Oracle package 대비 개발 허용 범위 내에 있다.
Topology F1만 높고 dynamics가 실패하면 완료가 아니다.

\section{TD13. Target runtime engine, final Gaussian transport와 rendering}

\textbf{목표.}
Versioned object package와 static GS에서 정확한 14-module dependency로 한 frame을 실행하고,
최종 total deformation을 Gaussian mean/covariance/frame에 한 번 transport한다.

\textbf{선행 조건.}
TD04, TD05, TD07, TD08, TD10-core, TD11과 TD12의 core standalone modules/contracts.
H/F/G를 채택한 configuration에서만 각각 TD06/TD09/TD10-G optional artifact를 추가로 요구한다.

\subsection*{통합 Runtime API}

\begin{lstlisting}[language=Python]
class WindDynamicsEngine:
    def setup(
        self,
        gs: CanonicalGaussianAsset,
        setup_metadata: SetupMetadata,
    ) -> ObjectPackage: ...

    def initialize_state(
        self,
        package: ObjectPackage,
    ) -> RuntimeState: ...

    def step(
        self,
        package: ObjectPackage,
        state: RuntimeState,
        wind: WindQuery,
        handle_motion: HandleMotion,
        dt: float,
    ) -> tuple[FrameOutput, RuntimeState]: ...
\end{lstlisting}

\subsection*{Exact typed frame trace}

\begin{lstlisting}
PreparedFrameState
 -> SurfaceState
 -> AnchorAeroForces
 -> ReducedAeroLoad
 -> GlobalPrediction
 -> SelectorDecision
 -> PatchForceBatch
 -> AssemblyResult
 -> ComplementForce
 -> LocalPrediction
 -> CorrectedDynamicState
 -> FinalAnchorState
 -> DeformedGaussianAsset
 -> RuntimeState / FrameDiagnostics / FrameOutput
\end{lstlisting}

\begin{itemize}
  \item[\todo] Debug mode에서 위 trace와 shape/device/version을 frame별 저장한다.
  \item[\todo] \code{AnchorAeroForces}는 V0 full-anchor analytic evaluation을 뜻하고, 조건부 H 사용 시 sample/gather 결과와 full-validation 결과를 별도 lineage로 기록한다.
  \item[\todo] \code{SelectorDecision}은 fixed-\(K\) Oracle/cached-compliance selector를 포함하고, full-counter는 offline/low-rate analysis, G는 optional backend로만 표기한다.
  \item[\todo] \code{PatchForceBatch}는 analytic/base channel을 항상 포함하며 F defect channel은 model이 채택된 configuration에서만 존재한다.
  \item[\todo] Production mode에서 필요한 diagnostic만 비동기/저비용으로 기록한다.
  \item[\todo] Target setup/runtime이 explicit target triangle-mesh reconstruction, MPM/FEM/CFD/LBM code path를 호출하지 않는지 trace test한다.
  \item[\todo] Inactive expert가 실제 skip되는지 kernel/call counter로 검사한다.
  \item[\todo] Factorization, sample, binding, model cache lifetime과 invalidation을 검사한다.
  \item[\todo] NaN/Inf, invalid covariance, inverted area, excessive force/energy를 fail-safe diagnostic으로 잡는다.
\end{itemize}

\subsection*{Module 12: final anchor state}

\begin{itemize}
  \item[\todo] Corrected \(q,z\)를 rest anchor에 단 한 번 lift한다.
  \item[\todo] Final position/velocity, \code{SurfaceTangentMap3x2}, cross-product normal/area와 별도의 \code{FullShellMap3x3}를 계산한다.
  \item[\todo] Global map 후 Local map을 순차로 Gaussian에 적용하지 않는다.
  \item[\todo] Hard/compliant attachment error와 seam/leakage를 기록한다.
  \item[\todo] Rank/orientation degeneracy를 confidence/fallback에 전달한다.
\end{itemize}

\subsection*{Module 13: affine Gaussian transport와 rendering}

\begin{itemize}
  \item[\todo] 고정 anchor support의 affine MLS로 Gaussian mean을 transport한다.
  \item[\todo] Covariance에는 명시적으로 구성한 \code{FullShellMap3x3}만 사용해 \(F_{3\times3}\Sigma^0F_{3\times3}^\top\)를 계산하고, rank-2 tangent map에 cofactor를 적용하지 않는다.
  \item[\todo] Covariance SPD/eigenvalue/scale clamp를 구현한다.
  \item[\todo] Co-rotational appearance/view frame을 갱신한다.
  \item[\todo] Identity/translation/rotation/in-plane affine reproduction을 다시 검사한다.
  \item[\todo] Center-only, rigid-only, full-affine backend를 ablation으로 제공한다.
  \item[\todo] Invalid covariance, scale explosion, frame flip, coverage hole/overlap과 temporal popping/flicker를 기록한다.
  \item[\todo] Compatible GPU에서 standard Gaussian renderer를 연결한다.
\end{itemize}

\subsection*{Module 14: final state와 diagnostics}

\begin{itemize}
  \item[\todo] Corrected Global/Local state와 Active/Decay/Dormant state를 저장한다.
  \item[\todo] Hysteresis/recurrent memory의 lifetime을 갱신한다.
  \item[\todo] Module별 p50/p95 latency, kernel/call 수와 memory를 기록한다.
  \item[\todo] Force/torque/energy/conservation/corrector residual을 기록한다.
  \item[\todo] Long rollout에서 state/index/memory가 누적되거나 leak되지 않는지 검사한다.
\end{itemize}

\subsection*{Runtime profiling contract}

다음을 개별 timing scope로 보고한다.

\begin{enumerate}
  \item State/surface reconstruction과 wind query.
  \item Full-anchor aero force와 generalized projection; 채택된 경우에만 조건부 H gather/reduction.
  \item Global predictor.
  \item Cached-compliance selector feature/selection; low-rate/full-counter와 채택된 G는 별도 비용.
  \item Active Local gather/scatter; 채택된 경우에만 조건부 F inference.
  \item Patch-copy condensation과 unique-anchor scatter.
  \item Constraint operator update, Schur factor/dual solve와 primal backprojection.
  \item Local integrator.
  \item Global corrector.
  \item Gaussian transport.
  \item Rendering.
\end{enumerate}

Warm-up, GPU synchronization, batch 1/multi-object, p50/p95, peak memory와 kernel count를 모두 기록한다.
Interactive/real-time wording은 target hardware, batch, \(N_G\), physical budget과 renderer 포함 범위의
사전 정의한 end-to-end p95 기준을 실측으로 통과한 경우에만 사용한다.
그 전에는 measured latency만 보고한다.

\textbf{완료 기준.}
Oracle package와 predicted package 양쪽에서 동일 API가 실행되고,
14-module trace, no-explicit-target-triangle-reconstruction trace test, affine transport,
long rollout, p50/p95 profile과 rendering output이 재현된다.

\section{TD14. 비교 실험, generalization과 paper evidence}

\textbf{목표.}
각 method component가 실제로 필요한지 반증 중심으로 평가하고,
solver 정확도와 ``without explicit target triangle-mesh reconstruction'' end-to-end 성능을 공정한 track으로 분리한다.

\textbf{선행 조건.} 해당 component의 standalone artifact와 TD13 integrated runtime.

\subsection*{평가 Track}

\begin{longtable}{L{0.14\linewidth}L{0.34\linewidth}L{0.44\linewidth}}
\toprule
Track & 통제 & 검증 목적 \\
\midrule
A: Oracle-topology matched solver & 동일 anchor scaffold/full-space operator, mass/material/attachment, wind, timestep, precision와 hardware; 실제 DOF/rank 별도 보고 & Aero/Global/Local/selector/coupling과 채택된 H/F/G의 solver 차이. \\
B: Static-3DGS end-to-end & 각 방법의 원래 preprocessing/input privilege를 명시 & Topology distillation, explicit target triangle-mesh reconstruction 요구, 전체 quality--latency. \\
C: Real captured static GS & Static multi-view/GS와 별도 dynamic reference & 실제 transfer, trajectory/silhouette/frequency, visual plausibility와 failure. \\
\bottomrule
\end{longtable}

\subsection*{E0--E10 실행 체크}

\begin{longtable}{L{0.07\linewidth}L{0.31\linewidth}L{0.52\linewidth}}
\toprule
ID & 실험 & 완료 증거 \\
\midrule
\endfirsthead
\toprule
ID & 실험 & 완료 증거 \\
\midrule
\endhead
E0 & Deterministic contracts & PSD/symmetry/orthogonality/force owner/transport/decay/double-count tests. \\
E1 & Topology/operator distillation & TD12-R과 Oracle/proposed/kNN/FPS/mesh baseline의 static+downstream error. \\
E2 & Full-anchor aero와 조건부 H & Full-anchor를 기준으로 rest/random/FPS/RSH/cubature의 generalized load, wrench, sample, latency. \\
E3 & Global physics vs neural & Physics/hybrid/direct/recurrent의 rollout, frequency, damping, OOD, latency. \\
E4 & Local physics와 조건부 F & None/full/reduced physics, optional missing-force hybrid, direct-displacement의 local spectrum, decay, cost. \\
E5 & Fixed-\(K\) selector와 조건부 G & Never/Always, matched-\(K\) random/curvature/strain/cached-compliance/Oracle, low-rate/full-counter analysis와 optional learned G의 regret, actual skip, Pareto. \\
E6 & Complement/overlap & No complement, post projection, integrate-average, no torque, proposed의 seam/wrench/leakage. \\
E7 & Feedback/corrector & No/structural/aero/next-frame/one/multi corrector의 Global effect, error, latency. \\
E8 & Static-3DGS end-to-end & Matched direct baselines, runtime privilege, dynamics/rendering/latency/memory. \\
E9 & Real static-GS transfer & Marker/feature trajectory, silhouette, tip/frequency/damping, novel-view temporal quality. \\
E10 & Quality--latency/scaling & Rank/local/active/corrector/Gaussian sweep과 채택된 H sample/F/G의 Pareto와 break-even. \\
\bottomrule
\end{longtable}

Track A의 solver baseline 우선순위는 E1의 Oracle/proposed/kNN/FPS/explicit target-mesh
representation baseline을 대체하거나 삭제하지 않는다.

\subsection*{Generalization와 stress test}

\begin{itemize}
  \item[\todo] Unseen yaw/elevation, magnitude와 combined wind OOD.
  \item[\todo] Unseen gust frequency/chirp와 stationary/slow/fast/reversing/localized gust.
  \item[\todo] Unseen setup material/stiffness/damping/density, setup attachment layout과 geometry.
  \item[\todo] Unseen category와 independent GS reconstruction/corruption.
  \item[\todo] Zero-wind relaxation, abrupt start/stop, threshold crossing.
  \item[\todo] Constant strong wind, resonance chirp, traveling gust.
  \item[\todo] 30--60초 rollout과 timestep \(\times0.5,\times2\).
  \item[\todo] Strong OOD에서 graceful failure, covariance/state validity와 fallback.
\end{itemize}

\subsection*{개발 go/no-go}

\begin{itemize}
  \item[\todo] TD12-R에서 independent GS의 area/mass, 첫 8--16 modes, modal principal angle과 close-layer false edge가 허용 범위인지 조기에 판정한다.
  \item[\todo] Predicted topology downstream error가 Oracle보다 약 10--15\% 이상 추가 악화되지 않는지 확인한다.
  \item[\todo] Same-anchor sparse PD/XPBD/corotational solver 대비
  개발용 내부 기준인 dynamics-only 약 2배 이상의 same-quality speedup,
  동일 p95 latency의 velocity/trajectory/spectrum 개선 또는
  multi-object throughput 이득 중 하나를 요구한다.
  \item[\todo] Fixed-\(K\) adaptive Local이 Always-on Local 대비 dynamics error 약 5\% 내외 증가에서 Local 비용 약 2배 절감을 목표로 한다.
  \item[\todo] H는 full-anchor profile 병목과 generalized-load fidelity 개선, F는 physics-only Local 개선, G는 cached-compliance selector 대비 quality--latency 개선을 각각 보일 때만 유지한다.
  \item[\todo] Gaussian 수 증가와 multi-object 실행에서 dynamics 비용이 rendering resolution과 거의 분리되는지 확인한다.
  \item[\todo] MVP Global tip frequency 약 5\%, amplitude 약 10\% 이내를 개발 목표로 사용한다.
  \item[\todo] 4--10초 rollout normalized position error가 object scale 약 5\% 내외이며 divergence가 없는지 확인한다.
\end{itemize}

위 수치는 논문 약속이 아니라 개발 중 방향을 중단하거나 단순화하기 위한 임시 기준이다.

\textbf{완료 기준.}
E0--E10 결과가 재현 가능한 manifest와 함께 생성되고,
아래 첫 논문의 두 algorithmic main claim(MC1/MC2), 하나의 SYS/supporting claim과
채택된 optional module마다
적어도 하나의 direct evidence, ablation, falsification result와 failure case가 있다.

\section{학습 curriculum과 loss implementation registry}

학습 가능한 component를 처음부터 joint training하지 않는다.
각 stage는 이전 deterministic/Oracle stage의 완료 증거를 요구한다.

\subsection{Stage 순서}

\begin{longtable}{L{0.06\linewidth}L{0.27\linewidth}L{0.49\linewidth}}
\toprule
Stage & 활성 component & 진입/종료 조건 \\
\midrule
\endfirsthead
\toprule
Stage & 활성 component & 진입/종료 조건 \\
\midrule
\endhead
0 & Deterministic contracts/structural preflight/transport & TD01 E0와 K00--K09 중 deterministic 해당 항목,
K20, K22--K23 및 K21 schema/synthetic fallback fixture를 통과한 뒤 Track A0에 진입한다.
Learned K21 calibration은 Track B0/E1, K24는 selector blind-spot stage에서 판정한다. \\
1 & Oracle topology + Global physics & TD05 Global-only start/stop/chirp/energy 완료. \\
1R & Early representation feasibility & Step 0에서 최소 schema/relation probe를 병렬 실행하고,
Step-0 gate 통과 뒤 5--10개 object의 formal TD12-R로 Gate A를 판정. \\
2 & Oracle Local에서 fixed-\(K\) selector/corrector로 전환 & TD07-A unit solver
\(\rightarrow\) TD11-A joint complement \(\rightarrow\) TD07-B same-p95 Gate B
\(\rightarrow\) TD10 cached-compliance/active-block skip/decay
\(\rightarrow\) TD11-B SYS ablation 순서를 재현함. \\
3 & 조건부 F always-on missing-force expert & TD08의 admissible generalized-to-nodal route를 고정한 뒤 TD09를 실행하며, K11 실패 시 F 없이 종료. \\
4 & 조건부 G benefit ranker/adaptive dispatch & Fixed-\(K\) cached-compliance/Oracle baseline과 positive-benefit label 뒤에만 시작; K12 실패 시 cached-compliance selector 유지. \\
5 & Full predicted topology/operator & TD12-R 통과와 Oracle package downstream sensitivity가 확정된 뒤 distill. \\
6 & Independent GS robustness & Reconstruction seed/density/hole/floater/object-disjoint split 평가. \\
7 & End-to-end transport/rendering & 모든 dynamics P0 통과 후 rendering quality/temporal loss를 연결. \\
\bottomrule
\end{longtable}

\subsection{Loss별 구현 체크}

\begin{itemize}
  \item[\todo] \textbf{Topology loss:} geometry/material-adjacency edge, hinge, orientation과 reconstruction reliability를 분리해 logging한다. User attachment/material 값 자체는 예측 loss에 넣지 않는다.
  \item[\todo] \textbf{Area/mass loss:} local 오차뿐 아니라 object total conservation을 포함한다.
  \item[\todo] \textbf{Operator/modal loss:} PSD construction, eigenvalue/mode shape, downstream rollout을 함께 본다.
  \item[\todo] \textbf{조건부 H cubature loss:} generalized load fidelity를 1차 목표로 두고 net force/torque, nonnegative/sparsity와 실제 latency를 보조 항으로 분리한다.
  \item[\todo] \textbf{조건부 F missing-force loss:} 선택한 nodal-teacher 또는 constrained-lift label route, channel-wise force, direction/magnitude, confidence와 inverse-dynamics residual을 기록한다.
  \item[\todo] \textbf{조건부 G gate loss:} positive-benefit regression/ranking, calibration과 fixed-\(K\) regret를 기록한다. Measured cost/budget loss는 profiler가 필요성을 보인 optional extension에서만 추가한다.
  \item[\todo] \textbf{Rollout loss:} one-step과 multi-step position/velocity/normal/curvature/spectrum을 분리한다.
  \item[\todo] \textbf{Energy/work loss:} unphysical positive work와 wind-stop decay violation을 penalize한다.
  \item[\todo] \textbf{Overlap/complement loss:} seam, Global leakage,
  free intrinsic wrench와 hard companion을 합친 sign-resolved wrench를 구분해 기록한다.
  \item[\todo] \textbf{Feedback loss:} predictor/corrected aero delta와 Global state error를 기록한다.
  \item[\todo] \textbf{Rendering loss:} dynamics가 통과한 뒤 PSNR/SSIM/LPIPS/temporal consistency를 추가한다.
  \item[\todo] Total loss weight와 schedule을 config/manifest에 저장하고 stage별 active term을 검증한다.
\end{itemize}

\subsection{Training 안정성과 재현성}

\begin{itemize}
  \item[\todo] Train/validation normalization statistic과 OOD clipping 정책을 저장한다.
  \item[\todo] Teacher-forced one-step와 closed-loop rollout batch 비율을 기록한다.
  \item[\todo] Gradient norm, NaN/Inf, energy violation과 rollout divergence를 monitor한다.
  \item[\todo] Learned component마다 최소 3개 seed를 실행할 수 있게 한다.
  \item[\todo] Best checkpoint selection metric을 paper test metric과 분리한다.
  \item[\todo] Test/OOD asset을 checkpoint selection에 사용하지 않는다.
\end{itemize}

\section{기존 방법 baseline 구현 registry}

Official code를 사용할 수 없거나 input contract를 맞추기 위해 재구현하면 논문명을 그대로 method label로 쓰지 않고
\code{-style} 또는 \code{matched reimplementation}이라고 표시한다.

\subsection{직접 또는 matched 정량 비교}

\small
\begin{longtable}{L{0.20\linewidth}L{0.18\linewidth}L{0.24\linewidth}L{0.18\linewidth}}
\toprule
방법 & 유사점 & 핵심 차이 & 필요한 비교 \\
\midrule
\endfirsthead
\toprule
방법 & 유사점 & 핵심 차이 & 필요한 비교 \\
\midrule
\endhead
High-resolution shell/MPM teacher & Wind-driven deformable surface & Privileged accuracy oracle/label source이며 runtime efficiency baseline 아님 & Accuracy upper bound, label/force balance. \\
Same-anchor sparse PD/XPBD & 동일 scaffold/material/load의 full-space physics; 실제 DOF 별도 보고 & Reduced Global/selective Local 없음 & MPM 대비 해상도 차이만이 아니라 matched-scaffold quality--latency 검증. \\
Wind Projection Basis-style & Directional wind를 modal load로 projection & Rest-state/uniform approximation, GS/Local 없음 & Current-state spatial aero와 Global dynamics E2/E3. \\
Gaussian Swaying-style + matched MPM & Current Gaussian normal/area/relative wind aero & Full MPM, Global ROM/gate 없음 & 같은 aero front-end의 accuracy/latency E8. \\
FreeForm-style mesh-free ROM & Point/GS 기반 reduced mechanics & Wind-specific topology/aero/Local gate 없음 & Oracle/predicted Global package와 ROM 비교. \\
MeshGraphNets-style rollout & Flag-like direct learned dynamics & Mesh/all-node neural, physics base 없음 & Physics 대 direct neural OOD/long stability E3. \\
Recurrent Global network & State/wind에서 next Global state & Restoring/damping을 data로 학습 & Memory가 있어도 physics Global보다 유리한지 E3. \\
Hybrid Global residual & Physics + learned generalized force & Global defect까지 학습 & Global network가 정말 필요한지 E3. \\
PGRD-style all-node residual & Physics + learned per-particle residual & All-node/all-time; spatial benefit gate/complement 없음 & Always-on residual 대 adaptive Local E4/E5. \\
Proposed Global-only & 동일 package/Global backbone & Local off & Local 순수 기여 E4. \\
Proposed Always-on Local physics & 동일 complementary Local system & 모든 patch active, 비용 상한 & Selector 품질 상한/비용 상한 E5. \\
Proposed cached-compliance selector & 동일 Local system, fixed \(K_A\) & Cached force와 prefactorized group operator를 쓰는 every-frame V0 selector & Curvature/strain/random 대비 core quality--latency E5. \\
Proposed full-counter residual & 동일 Local system, low-rate/offline & Extra Global solve와 force 재평가 비용 포함 & Teacher-free analysis ceiling E5. \\
Proposed Oracle selector & 동일 Local system, perfect future-benefit ranking & Runtime 불가능 & Patch selection upper bound E5. \\
Proposed optional F & Physics-only Local + missing-force expert & K11 통과 시에만 채택 & Missing force 순수 기여 E4/E8. \\
Proposed optional G & 동일 fixed \(K_A\), learned benefit ranker & K12 통과 시에만 채택 & Cached-compliance selector 대비 incremental Pareto E5/E10. \\
\bottomrule
\end{longtable}
\normalsize

\subsection{조건부 subset 비교}

\begin{longtable}{L{0.22\linewidth}L{0.29\linewidth}L{0.41\linewidth}}
\toprule
방법 & 공통화 가능한 subset & 비교 목적 \\
\midrule
DynamicTree & Tree/branch/leaf asset, external-force compact response & Static GS global modal-like response와 runtime mesh/Local 차이. \\
DiffWind & Prescribed spatial wind forward-simulation subset & Grid MPM/LBM 대 reduced one-way wind dynamics. \\
PhysSkin & 동일 external load/reduced deformation subset & Static GS의 point/edge reduced basis와 Gaussian lifting. \\
PhysGaussian & 동일 material/external force subset & MPM 기반 Gaussian dynamics 비용/품질. \\
GausSim & 동일 elastic object/input이 가능한 subset & Fixed hierarchy 대 error-triggered complementary Local. \\
NGFF & 동일 object-level/global-local force setting subset & Rigid-global/local stress 대 deformable Global modes와 sparse gate. \\
\bottomrule
\end{longtable}

공통 input/action/metric을 만들 수 없으면 정량 main table에 억지로 넣지 않고 capability/failure를 정성적으로만 비교한다.

\subsection{인용/설계 경계 중심}

\begin{itemize}
  \item Subspace Condensation: selective local full-space activation과 two-way coupling의 고전 경계.
  \item Trading Spaces: error/oracle 기반 adaptive enrichment의 경계.
  \item Learning Contact Corrections: learned local nonlinear correction의 경계.
  \item Fast Complementary Dynamics: Global/primary space의 complement constraint 경계.
  \item OmniAD: object-level directional aerodynamic SH representation 경계.
  \item Subspace/cloth two-level solvers: global low-frequency/local high-frequency 분할의 경계.
\end{itemize}

\section{Metric, profiling과 통계 계약}

\subsection{Metric registry}

\small
\begin{longtable}{L{0.17\linewidth}L{0.29\linewidth}L{0.46\linewidth}}
\toprule
범주 & Metric & 구현/보고 계약 \\
\midrule
\endfirsthead
\toprule
범주 & Metric & 구현/보고 계약 \\
\midrule
\endhead
Geometry & Object-scale normalized position/velocity RMSE & Frame와 sequence 평균을 구분하고 object scale 정의를 기록. \\
Gate B/C primary & Velocity RMSE 대 p95 dynamics latency & Position trajectory와 PSD/spectrum 대 p95를 공동 필수 보조 curve로 보고. \\
Surface & Normal angular, strain, curvature error & Invalid/degenerate sample 비율과 함께 보고. \\
Dynamics & Tip amplitude/phase, frequency, damping & Start/stop/sinusoidal/chirp별 time response와 PSD 저장. \\
Aero & Sample/generalized-force error & Full evaluator 기준 magnitude/direction/phase 분리. \\
Conservation & Net force/torque, COM/rotation drift & Internal/external channel과 torque center를 명시. \\
Energy & Work--energy residual, max growth, divergence & Zero wind, wind stop, activation event를 별도 보고. \\
Complement & Patch별 \(\Phi^\top M\Psi_r\), force leakage & Basis build와 runtime assembly 둘 다 검사하고 전역 basis mixing 여부를 기록. \\
Patch & Boundary position/velocity/force jump, active-block condition & Core--halo와 activation event에 대한 최대값/분포 및 cross-block group별 통계. \\
Selector/G & Regret, active ratio; optional G AUPRC/calibration & 동일 fixed \(K_A\)/latency 비교와 actual Local/F skip 포함. \\
Runtime & Module p50/p95, kernel/call, peak memory & Warm-up, GPU sync, batch/device/version 고정. \\
Topology & Edge/hinge F1, shortcut, attachment-transfer, area/mass, modal angle & User setup metadata를 고정하고 modal error--downstream rollout scatter를 asset/load별로 보고. \\
Gaussian & Invalid covariance, scale explosion, frame flip, coverage hole/overlap/popping & Count와 affected Gaussian/sequence 비율. \\
Rendering & PSNR/SSIM/LPIPS, temporal warp/flicker & Physics main table과 분리하되 end-to-end 결과에 포함. \\
Real data & Marker/feature trajectory, silhouette IoU & Ground-truth 한계를 명시하고 synthetic physics와 분리. \\
\bottomrule
\end{longtable}
\normalsize

Global rank와 Local coordinate는 support/실행 비용이 다르므로 동일 coordinate count를
공정성의 충분조건으로 사용하지 않는다. Rank/coordinate match는 진단용이며,
same-quality speedup, same-p95 velocity/trajectory/spectrum 또는 multi-object throughput의
Pareto non-dominated region을 primary 판정으로 사용한다.

\subsection{통계와 reporting}

\begin{itemize}
  \item[\todo] Frame을 독립 표본으로 취급하지 않고 먼저 sequence별 statistic을 계산한다.
  \item[\todo] Asset 단위 평균과 asset bootstrap 95\% confidence interval을 보고한다.
  \item[\todo] Learned baseline과 proposed model은 가능한 최소 3개 seed를 보고한다.
  \item[\todo] Mean, median, worst decile와 divergence/failure sequence 비율을 함께 보고한다.
  \item[\todo] p50/p95 latency는 동일 hardware/backend, warm-up, synchronization으로 측정한다.
  \item[\todo] Same-quality latency와 same-p95-latency quality 두 Pareto를 모두 만든다.
  \item[\todo] Preprocessing/setup/training/runtime cost를 분리한다.
  \item[\todo] Input privilege, explicit target triangle-mesh reconstruction/runtime grid, external model/data를 main comparison table에 명시한다.
\end{itemize}

\section{권장 artifact와 run directory layout}

\begin{lstlisting}
assets/
  canonical_gs/
  teacher_sources/

datasets/
  manifests/
  teacher_structure/
  teacher_fsi/
  gs_variants/
  projected_states/
  optional_F_missing_force_labels/
  selector_labels/
  optional_G_benefit_labels/
  feedback_labels/

packages/
  <asset_id>/<package_version>/
    package_manifest.json
    anchors.npz
    topology.npz
    operators.npz
    bases.npz
    patches.npz
    optional_H_aero_cubature.npz
    gaussian_binding.npz
    validation.json

models/
  topology/
  optional_F_missing_force/
  optional_G_gate/

runs/
  <run_id>/
    config.yaml
    run_manifest.json
    dataset_hash.txt
    package_hash.txt
    checkpoints/
    diagnostics/
    profiler/
    renders/
    report.md
\end{lstlisting}

\begin{itemize}
  \item[\done] 실제 repository split에 맞는 physical path와 artifact storage를 TD00의 \path{manifests/artifact_policy.json}에서 확정했다.
  \item[\done] Large dataset/checkpoint는 Git에 직접 넣지 않고 manifest와 retrieval rule만 version control하도록 TD00 policy를 고정했다.
  \item[\todo] Artifact schema version migration과 old cache invalidation test를 추가한다.
\end{itemize}

\section{Open Decisions와 결정 시한}

\begin{longtable}{L{0.06\linewidth}L{0.32\linewidth}L{0.23\linewidth}L{0.17\linewidth}}
\toprule
상태 & 결정 질문 & 권장 첫 선택 & 결정 시한 \\
\midrule
\endfirsthead
\toprule
상태 & 결정 질문 & 권장 첫 선택 & 결정 시한 \\
\midrule
\endhead
\blocked & Teacher-S solver/library와 material law는 무엇인가? & 검증 가능한 nonlinear shell/FEM; dense MPM은 matched baseline & TD02 전 \\
\blocked & Metric scale를 어떻게 제공하는가? & Known dimension/scale marker; 없으면 명시적 nondimensional preset & TD01 전 \\
\done & Canonical runtime integrator는 무엇인가? & Objective projective frozen target + implicit midpoint; backward Euler/rest-linear은 Step 0 ablation & Structural freeze contract에서 확정 \\
\done & Runtime material/attachment를 어디서 얻는가? & Setup에서 user metadata/layout을 고정하고 frame wind/handle만 변경 & TD00 contract에서 확정 \\
\done & Local force network target과 label lifting은 무엇인가? & Raw free-support nodal patch force가 기본; generalized lift는 channel conservation target을 갖춘 free/no-hard optional 경로, hard structural fallback은 no-go & TD08 contract에서 확정 \\
\done & Fixed-superset \(\Psi\)인가 dynamic active basis인가? & Patch-local fixed superset + active principal blocks; dynamic remap은 후속 & TD03/TD07 contract에서 확정 \\
\done & Surface map 차원과 normal/covariance 정의는 무엇인가? & 3\(\times\)2 tangent cross product와 별도 3\(\times\)3 full shell map & TD01/TD04 contract에서 확정 \\
\done & Global predictor가 cross coupling을 소비하는가? & Predictor가 소비하고 corrector는 updated-minus-consumed delta만 소비 & TD00 contract에서 확정 \\
\blocked & 조건부 F missing-force expert backbone은 무엇인가? & K11을 시험할 때만 small co-rotated patch MLP/graph model, memory off & TD09 착수 전 \\
\blocked & 조건부 G backbone은 무엇인가? & Fixed-\(K\) cached-compliance baseline 뒤 cheap pooled ranker; predicted cost는 profiler 이후 & TD10 optional G 착수 전 \\
\blocked & Topology distillation backbone은 무엇인가? & Candidate-edge classifier + positive operator construction & TD12 전 \\
\blocked & 조건부 H cubature solver/constraint tolerance는 무엇인가? & TD05에서 aero 병목 확인 후 nonnegative sparse fit; generalized load primary, wrench auxiliary & TD06 optional H 착수 전 \\
\blocked & Teacher-FSI를 첫 논문에 넣을 것인가? & Optional small subset; Teacher-S claim과 분리 & TD02 dataset freeze 전 \\
\blocked & 논문용 Gaussian renderer/GPU는 무엇인가? & Compatible CUDA backend, CC 6.1 blocker와 분리 & TD13 전 \\
\bottomrule
\end{longtable}

\section{즉시 실행할 첫 작업 순서}

아래 순서는 learned component를 붙이기 전에 연구 방향을 가장 싸게 반증하는 경로다.

\begin{enumerate}
  \item \textbf{TD00--TD01:} schema, force ledger, E0 fixtures와 test command를 만든다.
  \item \textbf{Step 0 structural freeze:} 최소 Oracle strip에서 typed \(\mathcal Q_{\mathrm{op}}\),
  rest-linear/projective operator와 target ledger, objective normal-curvature closure,
  implicit midpoint, fixed \(T_s,\Phi_s\), support-restricted per-RHS route와
  free/hard/compliant post-step reaction--\(r_{\mathrm{ROM}}\)을 순서대로 고정한다.
  \item \textbf{병렬 최소 조기 kill:} Step 0 fixture에 필요한 한쪽 고정 strip teacher와
  formal TD12-R 준비용 최소 independent-GS schema/relation probe만 병렬 생성한다.
  대규모 data/package 학습과 H/F/G는 보류한다.
  \item \textbf{TD02--TD03:} Step 0 통과 뒤 full teacher/data와
  Oracle anchor/operator/\(\Phi,\Psi\)/binding package를 확장한다.
  \item \textbf{TD04--TD05:} Full-anchor current-surface aero + Global-only ROM을
  same-anchor sparse solver와 비교해 K13 sub-gate의 개발용 약 2배 same-quality speedup,
  same-p95 velocity/trajectory/spectrum 또는 multi-object throughput을 먼저 판정한다.
  \item \textbf{TD07-A:} Patch-local fixed-superset operator와 single non-overlap Local unit solver를 만든다.
  \item \textbf{TD11-A:} Oracle/synthetic force로 unique-anchor small-Schur joint assembly와 Global complement를 고정한다.
  \item \textbf{TD07-B:} TD11-A force를 쓰는 Oracle coupled Local이 same-p95 Global rank sweep보다 나은지 Gate B를 판정한다.
  \item \textbf{TD10-core:} Cached-compliance selector와 Active/Decay/Dormant state, 실제 active/solve-set 비용을 학습 없이 완성한다.
  \item \textbf{TD11-B:} Gate B 뒤 corrected-aero/structural-delta one-corrector를 SYS ablation으로 검증한다.
  \item \textbf{TD08:} Nodal-breakdown 기본 label, auxiliary conservation target이 있는
  free/no-hard optional lift와 positive future-benefit/feedback label을 고정한다.
  \item \textbf{조건부 TD09/TD10-G:} F가 physics-only Local을 개선할 때만 F를, learned benefit이 cached-compliance selector를 개선할 때만 G를 채택한다.
  \item \textbf{조건부 TD06:} TD05 profiler에서 full-anchor aero가 실제 병목일 때만 H representation probe를 실행한다.
  \item \textbf{TD12--TD13:} TD12-R 통과 후 full predicted package를 학습하고 independent GS runtime을 연결한다.
  \item \textbf{TD14:} 각 단계 artifact를 E0--E10 표/그림으로 승격한다.
\end{enumerate}

\begin{tcolorbox}[colback=LightGreen,colframe=DoneGreen,title=첫 논문 최소 구현선,fonttitle=\bfseries]
Attached single-layer flag/cloth strip/paper/leaf, prescribed spatial wind, Teacher-S 중심,
independently reconstructed static 3DGS와 fixed metric/material/attachment metadata에서
explicit target triangle-mesh reconstruction과 runtime full-order continuum solve 없이 만드는 topology-distilled reduced physical scaffold,
full-anchor current-state aero와 Global physics, fixed-\(K\) cached-compliance selector가 선택하는 complementary Local physics,
conservative overlap/complement, corrected geometry-to-load feedback,
affine Gaussian transport와 standard rendering까지다.
첫 논문은 두 algorithmic main claim과 이를 완결하는 system axis로 정리한다.
(1) setup-conditioned, topology-distilled reduced physical scaffold,
(2) Global ROM 너머의 고주파 wind response를 선택적으로 복원하는 error-triggered complementary Local dynamics,
(SYS) coherent directly renderable animation을 위한 closed-loop geometry-to-load update와 affine Gaussian transport다.
Error-triggered의 최종 사용 여부는 learned G가 아니라 앞의 error-relevance evidence gate로 판정한다.
H/F/G는 각각 K14/K11/K12를 통과한 경우에만 supporting contribution 또는 variant로 추가하며 최소 구현선의 필수 요소가 아니다.
Full tree, self-contact, tearing, free flight, appearance residual과 arbitrary cross-family material generalization은 후속 범위다.
\end{tcolorbox}

\section{완료 정의와 최종 handoff}

각 milestone을 \done으로 바꾸려면 다음 다섯 증거가 모두 있어야 한다.

\begin{enumerate}
  \item Reusable implementation과 versioned interface.
  \item Reproduction command/config/seed/device가 있는 experiment README.
  \item Positive test와 의도적으로 실패하는 negative test.
  \item 정량 report와 최소 한 개의 visual/trace diagnostic.
  \item P0/go-no-go 결과, 실패 사례와 다음 decision을 기록한 session note.
\end{enumerate}

다음은 완료로 보지 않는다.

\begin{itemize}
  \item 조건부 F/G network loss만 감소하고 closed-loop rollout이 검증되지 않은 경우.
  \item 조건부 G score만 높고 Local/F kernel 또는 wall-clock가 줄지 않은 경우.
  \item Topology F1만 높고 downstream dynamics가 실패한 경우.
  \item Rendering이 그럴듯하지만 force, frequency, damping, energy가 맞지 않는 경우.
  \item Patch 결과를 시각적으로 평균해 seam을 숨긴 경우.
  \item Teacher privilege나 explicit target triangle-mesh reconstruction 사용을 manifest에서 숨긴 경우.
  \item Baseline을 original task/input와 다르게 구현하고 차이를 보고하지 않은 경우.
\end{itemize}

\section*{Notes}

\begin{itemize}
  \item 새 체크리스트는 이전 implementation checklist의 formatting과 검증 원칙만 계승한다.
  \item 기존 완료 상태, mesh-proxy extraction 중심 pipeline, mandatory SH residual은 계승하지 않는다.
  \item Global과 Local 모두 최종 motion은 physical integrator가 계산한다. 조건부 dynamics network F/G는 missing force 또는 positive future benefit만 근사한다.
  \item Fixed-\(K\) selector와 조건부 G는 새 Local excitation 및 채택된 F evaluation을 선택할 뿐 \(M,C,K\)나 structural base를 끄지 않는다.
  \item RSH는 공력 방향 표현 ablation이며 paper headline이나 main runtime contract가 아니다.
  \item 가장 먼저 반증할 질문은 \emph{same-anchor sparse solver보다 의미 있는 이득이 있는가},
  \emph{independent GS에서 low-mode scaffold를 복원할 수 있는가},
  \emph{Global rank 확대보다 Local이 필요한가}다.
  Physics-only Local보다 F가 필요한지는 그 다음의 조건부 질문이다.
\end{itemize}

\end{document}
